% Generated mechanically from the frozen criteria.tex target.
% Do not edit this generated file; edit the bound locale source instead.

% OK, start here.
%

\setcounter{chapter}{96}
\chapter{可表示性判据}


\phantomsection
\label{criteria-section-phantom}





\section{引言}
\label{criteria-section-introduction}

\noindent
本章的目的是寻找一些判据，用以保证赋予 fppf 拓扑的概形范畴上的
群胚叠是代数叠。历史上，这往往涉及证明某些函子可表示；参见
Grothendieck 的讲义
\cite{Gr-I},
\cite{Gr-II},
\cite{Gr-III},
\cite{Gr-IV},
\cite{Gr-V}，以及
\cite{Gr-VI}.
本章的标题由此而来。本章素材的另一个重要来源是 Artin 的工作；参见
\cite{ArtinI},
\cite{ArtinII},
\cite{Artin-Theorem-Representability},
\cite{Artin-Construction-Techniques},
\cite{Artin-Algebraic-Spaces},
\cite{Artin-Algebraic-Approximation},
\cite{Artin-Implicit-Function}，以及
\cite{ArtinVersal}.

\medskip\noindent
% P11-E0165
本章中的某些记号、约定和术语颇为别扭，经验更丰富的读者甚至可能觉得
它们前后颠倒。这是有意为之。解释见《Quot 空间与 Hilbert 空间》第
\ref{quot-section-conventions} 节。



\section{约定}
\label{criteria-section-conventions}

\noindent
本章采用的约定与《代数叠》一章中的约定相同；参见《代数叠》第
\ref{algebraic-section-conventions} 节。




\section{已有结果}
\label{criteria-section-done-so-far}

\noindent
代数空间情形下与本章对应的是题为《自举》的一章；参见《自举》第
\ref{bootstrap-section-introduction} 节。该章已经包含一些可表示性结果。
此外，其中处理的一部分预备材料，我们已经在《代数叠》一章中建立。
具体如下：
\begin{enumerate}
\item 我们在《自举》第
\ref{bootstrap-section-morphism-representable-by-spaces} 节讨论可由代数空间
表示的预层态射。在《代数叠》第
\ref{algebraic-section-morphisms-representable-by-algebraic-spaces}
节，我们讨论群胚纤维化范畴的 $1$-态射可由代数空间表示这一概念。
\item 我们在《自举》第
\ref{bootstrap-section-representable-by-spaces-properties} 节讨论可由代数空间
表示的预层态射的性质。在《代数叠》第
\ref{algebraic-section-representable-properties}
节，我们讨论可由代数空间表示的群胚纤维化范畴的 $1$-态射的性质。
\item 我们已经证明：若 $F$ 是一个层，其对角线可由代数空间表示，
并且它有一个来自代数空间的 \'etale 覆盖，则 $F$ 是代数空间；参见
《自举》，定理 \ref{bootstrap-theorem-bootstrap}。
（这是下一项结果的一个较弱版本。）
\item
\label{criteria-item-bootstrap-final}
我们已经证明：若 $F$ 是一个层，并且存在代数空间 $U$ 以及态射
$U \to F$，此态射可由代数空间表示、满、平坦且局部有限表示，
则 $F$ 是代数空间；参见《自举》，定理
\ref{bootstrap-theorem-final-bootstrap}。
\item 我们还证明了 (\ref{criteria-item-bootstrap-final}) 对代数叠的“光滑”类似物：
% P11-E0166
若 $\mathcal{X}$ 是 $(\Sch/S)_{fppf}$ 上的群胚叠，并且存在
$(\Sch/S)_{fppf}$ 上可由代数空间表示的群胚叠 $\mathcal{U}$，以及
$1$-态射 $u : \mathcal{U} \to \mathcal{X}$，此态射可由代数空间表示、
满且光滑，则 $\mathcal{X}$ 是代数叠；参见《代数叠》，引理
\ref{algebraic-lemma-smooth-surjective-morphism-implies-algebraic}.
\end{enumerate}
我们现在的第一项任务，是证明 (\ref{criteria-item-bootstrap-final}) 对一般代数叠的
类似物，即定理 \ref{criteria-theorem-bootstrap}。



\section{群胚叠的态射}
\label{criteria-section-1-morphisms}

\noindent
本节为预备内容，初读时可略过。

\begin{lemma}
\label{criteria-lemma-etale-permanence}
设 $\mathcal{X} \to \mathcal{Y} \to \mathcal{Z}$ 是
$(\Sch/S)_{fppf}$ 上群胚纤维化范畴的 $1$-态射。
若 $\mathcal{X} \to \mathcal{Z}$ 与 $\mathcal{Y} \to \mathcal{Z}$
可由代数空间表示且为 \'etale，则
$\mathcal{X} \to \mathcal{Y}$ 也具有这些性质。
\end{lemma}

\begin{proof}
设 $\mathcal{U}$ 是 $S$ 上可表示的群胚纤维化范畴，并设
$f : \mathcal{U} \to \mathcal{Y}$ 是 $1$-态射。我们必须证明
$\mathcal{X} \times_\mathcal{Y} \mathcal{U}$ 可由代数空间表示，
且在 $\mathcal{U}$ 上为 \'etale。考虑复合
$h : \mathcal{U} \to \mathcal{Z}$。则
$$
\mathcal{X} \times_\mathcal{Z} \mathcal{U}
\longrightarrow
\mathcal{Y} \times_\mathcal{Z} \mathcal{U}
$$
% P11-E0167
是两个群胚纤维化范畴之间的 $1$-态射；这两个范畴都可由代数空间表示，
且在 $\mathcal{U}$ 上均为 \'etale。因此，由空间的性质，引理
\ref{spaces-properties-lemma-etale-permanence}，该态射由代数空间的
\'etale 态射表示。最后，态射 $f$ 诱导态射
$\mathcal{U} \to \mathcal{Y} \times_\mathcal{Z} \mathcal{U}$，而且有
$$
\mathcal{X} \times_\mathcal{Y} \mathcal{U} =
(\mathcal{X} \times_\mathcal{Z} \mathcal{U})
\times_{(\mathcal{Y} \times_\mathcal{Z} \mathcal{U})}
\mathcal{U}.
$$
\end{proof}

\begin{lemma}
\label{criteria-lemma-stack-in-setoids-descent}
设 $\mathcal{X}, \mathcal{Y}, \mathcal{Z}$ 是
$(\Sch/S)_{fppf}$ 上的群胚叠，并设
$\mathcal{X} \to \mathcal{Y}$ 与 $\mathcal{Z} \to \mathcal{Y}$
是 $1$-态射。若
\begin{enumerate}
\item $\mathcal{Y}$、$\mathcal{Z}$ 分别可由 $S$ 上的代数空间
$Y$、$Z$ 表示；
% P11-E0168
\item 相联系的代数空间态射 $Y \to Z$ 满、平坦且局部有限表示；
\item $\mathcal{Y} \times_\mathcal{Z} \mathcal{X}$ 是集合胚叠；
\end{enumerate}
则 $\mathcal{X}$ 是集合胚叠。
\end{lemma}

\begin{proof}
这是《叠》引理 \ref{stacks-lemma-stack-in-setoids-descent} 的特例。
\end{proof}

\noindent
下一个引理是《代数叠》引理
\ref{algebraic-lemma-smooth-surjective-morphism-implies-algebraic}
的类似物，并将被更强的定理 \ref{criteria-theorem-bootstrap} 所取代。

\begin{lemma}
\label{criteria-lemma-flat-finite-presentation-surjective-diagonal}
设 $S$ 是概形，且 $u : \mathcal{U} \to \mathcal{X}$ 是
$(\Sch/S)_{fppf}$ 上群胚叠的 $1$-态射。若
\begin{enumerate}
\item $\mathcal{U}$ 可由代数空间表示；
\item $u$ 可由代数空间表示、满、平坦且局部有限表示；
\end{enumerate}
则
% P11-E0169
$\Delta : \mathcal{X} \to \mathcal{X} \times \mathcal{X}$
可由代数空间表示。
\end{lemma}

\begin{proof}
给定 $S$ 上的两个概形 $T_1$、$T_2$，
% P11-E0170
把相联系的可表示纤维范畴记为
$\mathcal{T}_i = (\Sch/T_i)_{fppf}$。设给定 $1$-态射
$f_i : \mathcal{T}_i \to \mathcal{X}$。由《代数叠》引理
\ref{algebraic-lemma-representable-diagonal}，只需证明 $2$-纤维积
$\mathcal{T}_1 \times_\mathcal{X} \mathcal{T}_2$
可由代数空间表示。由《叠》引理
\ref{stacks-lemma-2-fibre-product-stacks-in-setoids-over-stack-in-groupoids}
，它无论如何都是集合胚叠。因此，
$\mathcal{T}_1 \times_\mathcal{X} \mathcal{T}_2$ 对应于
$(\Sch/S)_{fppf}$ 上的某个层 $F$；参见《叠》引理
\ref{stacks-lemma-stack-in-setoids-characterize}。
设 $U$ 是表示 $\mathcal{U}$ 的代数空间。由假设，
$$
\mathcal{T}_i' = \mathcal{U} \times_{u, \mathcal{X}, f_i} \mathcal{T}_i
$$
可由 $S$ 上的代数空间 $T'_i$ 表示。因此，
$\mathcal{T}_1' \times_\mathcal{U} \mathcal{T}_2'$ 可由代数空间
$T'_1 \times_U T'_2$ 表示。考虑交换图
$$
\xymatrix{
&
\mathcal{T}_1 \times_{\mathcal X} \mathcal{T}_2 \ar[rr]\ar'[d][dd] & &
\mathcal{T}_1 \ar[dd] \\
\mathcal{T}_1' \times_\mathcal{U} \mathcal{T}_2' \ar[ur]\ar[rr]\ar[dd] & &
\mathcal{T}_1' \ar[ur]\ar[dd] \\
&
\mathcal{T}_2 \ar'[r][rr] & &
\mathcal X \\
\mathcal{T}_2' \ar[rr]\ar[ur] & &
\mathcal{U} \ar[ur] }
$$
在此图中，下方、右侧、后侧和前侧的方形都是 $2$-纤维积。于是形式论证表明，
$\mathcal{T}_1' \times_\mathcal{U} \mathcal{T}_2' \to
\mathcal{T}_1 \times_{\mathcal X} \mathcal{T}_2$
是 $\mathcal{U} \to \mathcal{X}$ 的“基变换”；更确切地，图
$$
\xymatrix{
\mathcal{T}_1' \times_\mathcal{U} \mathcal{T}_2' \ar[d] \ar[r] &
\mathcal{U} \ar[d] \\
\mathcal{T}_1 \times_{\mathcal X} \mathcal{T}_2 \ar[r] &
\mathcal{X}
}
$$
是一个 $2$-纤维积方形。因此，$T'_1 \times_U T'_2 \to F$
可由代数空间表示、平坦、局部有限表示且为满；参见《代数叠》引理
\ref{algebraic-lemma-map-fibred-setoids-representable-algebraic-spaces},
\ref{algebraic-lemma-base-change-representable-by-spaces},
\ref{algebraic-lemma-map-fibred-setoids-property} 和
\ref{algebraic-lemma-base-change-representable-transformations-property}.
所以，由《自举》定理 \ref{bootstrap-theorem-final-bootstrap}，
$F$ 是代数空间，结论得证。
\end{proof}

\begin{lemma}
\label{criteria-lemma-second-diagonal}
设 $\mathcal{X}$ 是 $(\Sch/S)_{fppf}$ 上的群胚纤维化范畴。
下列条件等价：
\begin{enumerate}
\item $\Delta_\Delta : \mathcal{X} \to
\mathcal{X} \times_{\mathcal{X} \times \mathcal{X}} \mathcal{X}$
可由代数空间表示；
\item 对每个 $1$-态射 $\mathcal{V} \to \mathcal{X} \times \mathcal{X}$，
若 $\mathcal{V}$（由概形）可表示，则纤维积
$\mathcal{Y} =
\mathcal{X} \times_{\Delta, \mathcal{X} \times \mathcal{X}} \mathcal{V}$
的对角态射可由代数空间表示。
\end{enumerate}
\end{lemma}

\begin{proof}
尽管这略费思量，但完全是形式的。回顾
$\mathcal{X} \times_{\mathcal{X} \times \mathcal{X}} \mathcal{X} =
\mathcal{I}_\mathcal{X}$ 是 $\mathcal{X}$ 的惯性，并且
$\Delta_\Delta$ 是 $\mathcal{I}_\mathcal{X}$ 的恒等截面；参见
《范畴》第 \ref{categories-section-inertia} 节。因此，条件 (1) 的含义如下：
给定概形 $V$、$\mathcal{X}$ 在 $V$ 上的对象 $x$ 以及
$\mathcal{X}_V$ 的态射 $\alpha : x \to x$，条件
“$\alpha = \text{id}_x$”在 $V$ 上定义一个代数空间。
（换言之，存在代数空间的单态射 $W \to V$，使概形态射
$f : T \to V$ 通过 $W$ 分解，当且仅当
$f^*\alpha = \text{id}_{f^*x}$。）

\medskip\noindent
另一方面，设 $V$ 是概形，$x,y$ 是 $\mathcal{X}$ 在 $V$ 上的对象。
于是 $(x,y)$ 定义态射
$\mathcal{V} = (\Sch/V)_{fppf} \to \mathcal{X} \times \mathcal{X}$.
再设 $h : V' \to V$ 是概形态射，并设
$\alpha : h^*x \to h^*y$ 和 $\beta : h^*x \to h^*y$ 是
$\mathcal{X}_{V'}$ 的态射。于是 $(\alpha, \beta)$ 定义态射
% P11-E0171: frozen source appears to use V rather than V' below and again in the following sentence; preserved canonically.
$\mathcal{V}' = (\Sch/V)_{fppf} \to \mathcal{Y} \times \mathcal{Y}$.
条件 (2) 现在断言（对上述任意选择），条件“$\alpha = \beta$”
在 $V$ 上定义一个代数空间。

\medskip\noindent
为证明等价性，给定 (2) 中的 $(\alpha,\beta)$，由 (1) 可知条件
“$\alpha^{-1} \circ \beta = \text{id}_{h^*x}$”定义一个代数空间。
取 $h = \text{id}_V$ 且 $\beta = \text{id}_x$，即得
(2) $\Rightarrow$ (1)。
\end{proof}











\section{在对象上保极限}
\label{criteria-section-limit-preserving}

\noindent
设 $S$ 是概形，$p : \mathcal{X} \to \mathcal{Y}$ 是
$(\Sch/S)_{fppf}$ 上群胚纤维化范畴的 $1$-态射。如果下列条件成立，
则称 $p$ {\it 在对象上保极限}：给定由下列各项组成的任意数据
\begin{enumerate}
\item 仿射概形 $U = \lim_{i \in I} U_i$，它写成 $S$ 上仿射概形
$U_i$ 的有向极限；
\item 对某个 $i$，$\mathcal{Y}$ 在 $U_i$ 上的对象 $y_i$；
\item $\mathcal{X}$ 在 $U$ 上的对象 $x$；以及
\item 同构 $\gamma : p(x) \to y_i|_U$，
\end{enumerate}
则存在 $i' \geq i$、$\mathcal{X}$ 在 $U_{i'}$ 上的对象 $x_{i'}$、同构
$\beta : x_{i'}|_U \to x$ 以及同构
$\gamma_{i'} : p(x_{i'}) \to y_i|_{U_{i'}}$，使得
\begin{equation}
\label{criteria-equation-limit-preserving}
\vcenter{
\xymatrix{
p(x_{i'}|_U) \ar[d]_{p(\beta)} \ar[rr]_{\gamma_{i'}|_U} & &
(y_i|_{U_{i'}})|_U \ar@{=}[d] \\
p(x) \ar[rr]^\gamma & & y_i|_U
}
}
\end{equation}
交换。在这种情形下，我们称“$(i',x_{i'},\beta,\gamma_{i'})$
是由数据 (1)、(2)、(3)、(4) 所提出问题的一个{\it 解}”。
这一定义的动机来自《空间的极限》引理
\ref{spaces-limits-lemma-characterize-relative-limit-preserving}。

\begin{lemma}
\label{criteria-lemma-base-change-limit-preserving}
设 $p : \mathcal{X} \to \mathcal{Y}$ 和
$q : \mathcal{Z} \to \mathcal{Y}$ 是 $(\Sch/S)_{fppf}$ 上
群胚纤维化范畴的 $1$-态射。如果
$p : \mathcal{X} \to \mathcal{Y}$ 在对象上保极限，则 $p$ 沿 $q$ 的基变换
$p' : \mathcal{X} \times_\mathcal{Y} \mathcal{Z} \to \mathcal{Z}$
也在对象上保极限。
\end{lemma}

\begin{proof}
这是形式的。设 $U = \lim_{i \in I} U_i$ 是 $S$ 上仿射概形
$U_i$ 的有向极限；对某个 $i$，设 $z_i$ 是 $\mathcal{Z}$ 在
$U_i$ 上的对象，$w$ 是 $\mathcal{X} \times_\mathcal{Y} \mathcal{Z}$
在 $U$ 上的对象，而
$\delta : p'(w) \to z_i|_U$ 是同构。可以写成
$w = (U, x, z, \alpha)$，其中 $x$ 是 $\mathcal{X}$ 在 $U$ 上的某个对象，
$z$ 是 $\mathcal{Z}$ 在 $U$ 上的对象，而
$\alpha : p(x) \to q(z)$ 是同构。注意 $p'(w)=z$，故
$\delta : z \to z_i|_U$。令 $y_i=q(z_i)$ 且
$\gamma = q(\delta) \circ \alpha : p(x) \to y_i|_U$。
由于 $p$ 在对象上保极限，存在 $i'\geq i$ 和 $\mathcal{X}$ 在
$U_{i'}$ 上的对象 $x_{i'}$，以及同构
$\beta : x_{i'}|_U \to x$ 和
$\gamma_{i'} : p(x_{i'}) \to y_i|_{U_{i'}}$，使
(\ref{criteria-equation-limit-preserving}) 交换。于是考虑
$\mathcal{X} \times_\mathcal{Y} \mathcal{Z}$ 在 $U_{i'}$ 上的对象
$w_{i'} = (U_{i'}, x_{i'}, z_i|_{U_{i'}}, \gamma_{i'})$，并定义同构
$$
w_{i'}|_U = (U, x_{i'}|_U, z_i|_U, \gamma_{i'}|_U)
\xrightarrow{(\beta, \delta^{-1})}
(U, x, z, \alpha) = w
$$
以及
$$
p'(w_{i'}) = z_i|_{U_{i'}} \xrightarrow{\text{id}} z_i|_{U_{i'}}.
$$
二者合在一起给出该问题的一个解。
\end{proof}

\begin{lemma}
\label{criteria-lemma-composition-limit-preserving}
设 $p : \mathcal{X} \to \mathcal{Y}$ 和
$q : \mathcal{Y} \to \mathcal{Z}$ 是 $(\Sch/S)_{fppf}$ 上
群胚纤维化范畴的 $1$-态射。如果 $p$ 和 $q$ 均在对象上保极限，
则复合 $q \circ p$ 也在对象上保极限。
\end{lemma}

\begin{proof}
这是形式的。设 $U = \lim_{i \in I} U_i$ 是 $S$ 上仿射概形
$U_i$ 的有向极限；对某个 $i$，设 $z_i$ 是 $\mathcal{Z}$ 在
$U_i$ 上的对象，$x$ 是 $\mathcal{X}$ 在 $U$ 上的对象，而
$\gamma : q(p(x)) \to z_i|_U$ 是同构。由于 $q$ 在对象上保极限，
存在 $i' \geq i$、$\mathcal{Y}$ 在 $U_{i'}$ 上的对象 $y_{i'}$、同构
$\beta : y_{i'}|_U \to p(x)$ 以及同构
$\gamma_{i'} : q(y_{i'}) \to z_i|_{U_{i'}}$，使
(\ref{criteria-equation-limit-preserving}) 交换。由于 $p$ 在对象上保极限，
存在 $i'' \geq i'$、$\mathcal{X}$ 在 $U_{i''}$ 上的对象 $x_{i''}$、同构
$\beta' : x_{i''}|_U \to x$ 以及同构
$\gamma'_{i''} : p(x_{i''}) \to y_{i'}|_{U_{i''}}$，使
(\ref{criteria-equation-limit-preserving}) 交换。所求解取为 $U_{i''}$ 上的
$x_{i''}$，并配以同构
$$
q(p(x_{i''})) \xrightarrow{q(\gamma'_{i''})}
q(y_{i'})|_{U_{i''}} \xrightarrow{\gamma_{i'}|_{U_{i''}}}
z_i|_{U_{i''}}
$$
和同构 $\beta' : x_{i''}|_U \to x$。我们略去
(\ref{criteria-equation-limit-preserving}) 交换性的验证。
\end{proof}

\begin{lemma}
\label{criteria-lemma-representable-by-spaces-limit-preserving}
设 $p : \mathcal{X} \to \mathcal{Y}$ 是 $(\Sch/S)_{fppf}$ 上
群胚纤维化范畴的 $1$-态射。如果 $p$ 可由代数空间表示，
则下列条件等价：
\begin{enumerate}
\item $p$ 在对象上保极限；
\item $p$ 局部有限表示（参见《代数叠》定义
\ref{algebraic-definition-relative-representable-property}）。
\end{enumerate}
\end{lemma}

\begin{proof}
假设 (2)。设 $U = \lim_{i \in I} U_i$ 是 $S$ 上仿射概形
$U_i$ 的有向极限；对某个 $i$，设 $y_i$ 是 $\mathcal{Y}$ 在
$U_i$ 上的对象，$x$ 是 $\mathcal{X}$ 在 $U$ 上的对象，而
$\gamma : p(x) \to y_i|_U$ 是同构。记 $X_{y_i}$ 为表示 $2$-纤维积
$$
(\Sch/U_i)_{fppf} \times_{y_i, \mathcal{Y}, p} \mathcal{X}.
$$
的 $U_i$ 上的代数空间。注意，
$\xi = (U, U \to U_i, x, \gamma^{-1})$ 定义此 $2$-纤维积在
$U$ 上的一个对象。由 $2$-Yoneda 引理，$\xi$ 对应于 $U_i$ 上的态射
$f_\xi : U \to X_{y_i}$。由《空间的极限》命题
\ref{spaces-limits-proposition-characterize-locally-finite-presentation}
，存在 $i'\geq i$ 和态射 $f_{i'} : U_{i'} \to X_{y_i}$，使
$f_\xi$ 是 $f_{i'}$ 与投影态射 $U \to U_{i'}$ 的复合。
$2$-Yoneda 引理还告诉我们，$f_{i'}$ 对应于上述 $2$-纤维积在
$U_{i'}$ 上的一个对象
$\xi_{i'} = (U_{i'}, U_{i'} \to U_i, x_{i'}, \alpha)$，其到 $U$ 的限制
恢复 $\xi$。特别地，我们得到同构
% P11-E0172: frozen source uses |U rather than |_U in the following two displays; preserved canonically.
$\gamma : x_{i'}|U \to x$。注意，$\alpha : y_i|_{U_{i'}} \to p(x_{i'})$。
因此，取 $x_{i'}$、同构
$\gamma : x_{i'}|U \to x$ 以及同构
$\beta = \alpha^{-1} : p(x_{i'}) \to y_i|_{U_{i'}}$，即为该问题的一个解。

\medskip\noindent
假设 (1)。选取概形 $T$ 和 $1$-态射
$y : (\Sch/T)_{fppf} \to \mathcal{Y}$。设 $X_y$ 是表示 $2$-纤维积
$(\Sch/T)_{fppf} \times_{y, \mathcal{Y}, p} \mathcal{X}$.
的 $T$ 上的代数空间。我们必须证明 $X_y \to T$ 局部有限表示。
为此，将使用《空间的极限》评注
\ref{spaces-limits-remark-limit-preserving} 中的判据。考虑仿射概形
$U = \lim_{i \in I} U_i$，它写成 $T$ 上仿射概形的有向极限。
任取 $i \in I$ 并令 $y_i = y|_{U_i}$。另以 $i'$ 表示 $I$ 中
大于或等于 $i$ 的元素。由 $2$-Yoneda 引理，$T$ 上的态射
$U \to X_y$ 与二元组 $(x,\alpha)$ 的同构类一一对应，其中 $x$ 是
$\mathcal{X}$ 在 $U$ 上的对象，而
$\alpha : y|_U \to p(x)$ 是同构。当然，在取逆的意义下，给出
$\alpha$ 等同于给出同构 $\gamma : p(x) \to y_i|_U$。
$T$ 上的态射 $U_{i'} \to X_y$ 情形类似。因此，(1) 保证典范映射
$$
\colim_{i' \geq i} X_y(U_{i'}) \longrightarrow X_y(U)
$$
在此情形下为满。由《空间的极限》引理
\ref{spaces-limits-lemma-surjection-is-enough}，$X_y \to T$ 局部有限表示。
\end{proof}

\begin{lemma}
\label{criteria-lemma-open-immersion-limit-preserving}
设 $p : \mathcal{X} \to \mathcal{Y}$ 是 $(\Sch/S)_{fppf}$ 上
群胚纤维化范畴的 $1$-态射。假设 $p$ 可由代数空间表示，且为开浸入。
则 $p$ 在对象上保极限。
\end{lemma}

\begin{proof}
这由引理 \ref{criteria-lemma-representable-by-spaces-limit-preserving}，并经由
《代数叠》引理
\ref{algebraic-lemma-representable-transformations-property-implication}
中的一般原理，再结合下述事实推出：代数空间的开浸入局部有限表示；
参见《空间的态射》引理
\ref{spaces-morphisms-lemma-open-immersion-locally-finite-presentation}。
\end{proof}

\noindent
设 $S$ 是概形。在下一个引理中，我们需要 $S$ 上代数空间 $X$ 的
{\it 大小}这一概念。确切地，给定基数 $\kappa$，我们称 $X$ 满足
$\text{size}(X) \leq \kappa$，当且仅当存在概形 $U$，使
$\text{size}(U) \leq \kappa$（参见《集合》第
\ref{sets-section-categories-schemes} 节），并且存在满的平展态射 $U \to X$。

\begin{lemma}
\label{criteria-lemma-check-representable-limit-preserving}
设 $S$ 是概形。对某个 $T \in \Ob((\Sch/S)_{fppf})$，令
$\kappa = \text{size}(T)$。设
$f : \mathcal{X} \to \mathcal{Y}$ 是 $(\Sch/S)_{fppf}$ 上
群胚纤维化范畴的 $1$-态射，并满足
\begin{enumerate}
\item $\mathcal{Y} \to (\Sch/S)_{fppf}$ 在对象上保极限；
\item 对 $S$ 上局部有限表示的仿射概形 $V$ 和
$y \in \Ob(\mathcal{Y}_V)$，纤维积
$(\Sch/V)_{fppf} \times_{y, \mathcal{Y}} \mathcal{X}$ 可由大小
$\leq \kappa$ 的代数空间表示\footnote{忽略集合论问题的读者可以删去大小条件。}；
\item $\mathcal{X}$ 和 $\mathcal{Y}$ 都是 Zariski 拓扑的叠。
\end{enumerate}
则 $f$ 可由代数空间表示。
\end{lemma}

\begin{proof}
设 $V$ 是 $S$ 上的概形且 $y \in \mathcal{Y}_V$。我们必须证明
$(\Sch/V)_{fppf} \times_{y, \mathcal{Y}} \mathcal{X}$ 可由代数空间表示。

\medskip\noindent
情形 I：$V$ 为仿射概形，并映入仿射开集
$\Spec(\Lambda) \subset S$。于是可写 $V = \lim V_i$，其中每个 $V_i$
都是 $\Spec(\Lambda)$ 上有限表示的仿射概形；参见《代数》引理
\ref{algebra-lemma-ring-colimit-fp}。由假设 (1)，对某个 $i$，$y$
来自 $V_i$ 上的对象 $y_i$。由假设 (3)，纤维积
$(\Sch/V_i)_{fppf} \times_{y_i, \mathcal{Y}} \mathcal{X}$ 可由代数空间
$Z_i$ 表示。于是
$(\Sch/V)_{fppf} \times_{y, \mathcal{Y}} \mathcal{X}$ 可由
% P11-E0173: frozen source writes Z below after defining Z_i; likely Z_i is intended, but the canonical formula is preserved.
$Z \times_{V_i} V$ 表示。

\medskip\noindent
情形 II：$V$ 为一般概形。选取仿射开覆盖
$V = \bigcup_{i \in I} V_i$，使每个 $V_i$ 都映入 $S$ 的一个仿射开集。
先断言
$\mathcal{Z} = (\Sch/V)_{fppf} \times_{y, \mathcal{Y}} \mathcal{X}$
是 Zariski 拓扑的集合胚叠。事实上，由《叠》引理
\ref{stacks-lemma-2-product-stacks-in-groupoids}，它是 Zariski 拓扑的
群胚叠。再设 $z$ 是 $\mathcal{Z}$ 在概形 $T$ 上的对象。
记 $g : T \to V$ 为与 $z$ 到 $(\Sch/V)_{fppf}$ 的投影相对应的态射。
考虑 Zariski 层
$\mathit{I} = \mathit{Isom}_{\mathcal{Z}}(z, z)$。由情形 I，
$\mathit{I}|_{g^{-1}(V_i)} = *$（单点层）。故
% P11-E0174: frozen source changes mathit I to mathcal I here without defining a new object; preserved canonically.
$\mathcal{I} = *$。所以 $\mathcal{Z}$ 纤维化于集合胚。为完成证明，
须证明 Zariski 层 $Z : T \mapsto \Ob(\mathcal{Z}_T)/\cong$ 是代数空间；
参见《代数叠》引理
\ref{algebraic-lemma-characterize-representable-by-space}。
存在映射 $p : Z \to V$（函子之间的变换），而由情形 I，
$Z_i = p^{-1}(V_i)$ 是代数空间。态射 $Z_i \to Z$ 可由开浸入表示，
且 $\coprod Z_i \to Z$ 在 Zariski 拓扑下为满。因此，由《自举》引理
\ref{bootstrap-lemma-glueing-sheaves}，$Z$ 是 fppf 拓扑的层。
于是可应用《空间》引理 \ref{spaces-lemma-glueing-algebraic-spaces}，
从而 $Z$ 是代数空间\footnote{为验证该引理中的集合论条件，论证如下。
首先选取开覆盖，使 $|I| \leq \text{size}(V)$。其次，选取大小
$\leq \max(\kappa, \text{size}(V))$ 的概形 $U_i$ 和满平展态射
$U_i \to Z_i$；由假设 (2) 和《集合》引理
\ref{sets-lemma-bound-size-fibre-product} 可以做到这一点（略去细节）。
于是《集合》引理 \ref{sets-lemma-what-is-in-it} 表明
$\coprod U_i$ 是 $(\Sch/S)_{fppf}$ 的对象。因此，由《空间》引理
\ref{spaces-lemma-coproduct-algebraic-spaces}，$\coprod Z_i$ 是代数空间。}。
\end{proof}

\begin{lemma}
\label{criteria-lemma-check-property-limit-preserving}
设 $S$ 是概形，$f : \mathcal{X} \to \mathcal{Y}$ 是
$(\Sch/S)_{fppf}$ 上群胚纤维化范畴的 $1$-态射。设 $\mathcal{P}$ 是
《代数叠》定义 \ref{algebraic-definition-relative-representable-property}
中所述的代数空间态射的性质。如果
\begin{enumerate}
\item $f$ 可由代数空间表示；
\item $\mathcal{Y} \to (\Sch/S)_{fppf}$ 在对象上保极限；
\item 对 $S$ 上局部有限表示的仿射概形 $V$ 和
$y \in \mathcal{Y}_V$，所得代数空间态射 $f_y : F_y \to V$
具有性质 $\mathcal{P}$；参见《代数叠》公式
(\ref{algebraic-equation-representable-by-algebraic-spaces})。
\end{enumerate}
则 $f$ 具有性质 $\mathcal{P}$。
\end{lemma}

\begin{proof}
设 $V$ 是 $S$ 上的概形且 $y \in \mathcal{Y}_V$。我们必须证明
$F_y \to V$ 具有性质 $\mathcal{P}$。由于 $\mathcal{P}$ 在基底上为
fppf 局部的，可以假设 $V$ 是映入仿射开集
$\Spec(\Lambda) \subset S$ 的仿射概形。于是可写
$V = \lim V_i$，其中每个 $V_i$ 都是 $\Spec(\Lambda)$ 上有限表示的
仿射概形；参见《代数》引理 \ref{algebra-lemma-ring-colimit-fp}。
由假设 (2)，对某个 $i$，$y$ 来自 $V_i$ 上的对象 $y_i$。
由假设 (3)，态射 $F_{y_i} \to V_i$ 具有性质 $\mathcal{P}$。
由于 $\mathcal{P}$ 在任意基变换下稳定，并且
$F_y = F_{y_i} \times_{V_i} V$，故 $F_y \to V$ 如愿具有性质
$\mathcal{P}$。
\end{proof}



\section{在对象上形式光滑}
\label{criteria-section-formally-smooth}

\noindent
设 $S$ 是概形，$p : \mathcal{X} \to \mathcal{Y}$ 是
$(\Sch/S)_{fppf}$ 上群胚纤维化范畴的 $1$-态射。如果下列条件成立，
则称 $p$ {\it 在对象上形式光滑}：给定由下列各项组成的任意数据
\begin{enumerate}
\item $S$ 上仿射概形的一阶加厚 $U \subset U'$；
\item $\mathcal{Y}$ 在 $U'$ 上的对象 $y'$；
\item $\mathcal{X}$ 在 $U$ 上的对象 $x$；以及
\item 同构 $\gamma : p(x) \to y'|_U$，
\end{enumerate}
则存在 $\mathcal{X}$ 在 $U'$ 上的对象 $x'$，并有同构
$\beta : x'|_U \to x$ 和同构 $\gamma' : p(x') \to y'$，使得
\begin{equation}
\label{criteria-equation-formally-smooth}
\vcenter{
\xymatrix{
p(x'|_U) \ar[d]_{p(\beta)} \ar[rr]_{\gamma'|_U} & &
y'|_U \ar@{=}[d] \\
p(x) \ar[rr]^\gamma & & y'|_U
}
}
\end{equation}
交换。在这种情形下，我们称“$(x',\beta,\gamma')$ 是由数据
(1)、(2)、(3)、(4) 所提出问题的一个{\it 解}”。

\begin{lemma}
\label{criteria-lemma-base-change-formally-smooth}
设 $p : \mathcal{X} \to \mathcal{Y}$ 和
$q : \mathcal{Z} \to \mathcal{Y}$ 是 $(\Sch/S)_{fppf}$ 上
群胚纤维化范畴的 $1$-态射。如果
$p : \mathcal{X} \to \mathcal{Y}$ 在对象上形式光滑，则 $p$ 沿 $q$ 的基变换
$p' : \mathcal{X} \times_\mathcal{Y} \mathcal{Z} \to \mathcal{Z}$
也在对象上形式光滑。
\end{lemma}

\begin{proof}
这是形式的。设 $U \subset U'$ 是 $S$ 上仿射概形的一阶加厚，
$z'$ 是 $\mathcal{Z}$ 在 $U'$ 上的对象，$w$ 是
$\mathcal{X} \times_\mathcal{Y} \mathcal{Z}$ 在 $U$ 上的对象，而
$\delta : p'(w) \to z'|_U$ 是同构。可以写
$w = (U, x, z, \alpha)$，其中 $x$ 是 $\mathcal{X}$ 在 $U$ 上的某个对象，
$z$ 是 $\mathcal{Z}$ 在 $U$ 上的对象，而
$\alpha : p(x) \to q(z)$ 是同构。注意 $p'(w)=z$，故
% P11-E0175: frozen source writes z|_U below although delta was defined with target z'|_U; preserved canonically.
$\delta : z \to z|_U$。令 $y'=q(z')$ 且
$\gamma = q(\delta) \circ \alpha : p(x) \to y'|_U$。
由于 $p$ 在对象上形式光滑，存在 $\mathcal{X}$ 在 $U'$ 上的对象 $x'$，
以及同构 $\beta : x'|_U \to x$ 和 $\gamma' : p(x') \to y'$，使
(\ref{criteria-equation-formally-smooth}) 交换。于是考虑
% P11-E0176: frozen source names the following lifted object w, while subsequent formulas use w'; preserved canonically.
$\mathcal{X} \times_\mathcal{Y} \mathcal{Z}$ 在 $U'$ 上的对象
$w = (U', x', z', \gamma')$，并定义同构
$$
w'|_U = (U, x'|_U, z'|_U, \gamma'|_U)
\xrightarrow{(\beta, \delta^{-1})}
(U, x, z, \alpha) = w
$$
以及
$$
p'(w') = z' \xrightarrow{\text{id}} z'.
$$
二者合在一起给出该问题的一个解。
\end{proof}

\begin{lemma}
\label{criteria-lemma-composition-formally-smooth}
设 $p : \mathcal{X} \to \mathcal{Y}$ 和
$q : \mathcal{Y} \to \mathcal{Z}$ 是 $(\Sch/S)_{fppf}$ 上
群胚纤维化范畴的 $1$-态射。如果 $p$ 和 $q$ 均在对象上形式光滑，
则复合 $q \circ p$ 也在对象上形式光滑。
\end{lemma}

\begin{proof}
这是形式的。设 $U \subset U'$ 是 $S$ 上仿射概形的一阶加厚，
$z'$ 是 $\mathcal{Z}$ 在 $U'$ 上的对象，$x$ 是 $\mathcal{X}$ 在
$U$ 上的对象，而 $\gamma : q(p(x)) \to z'|_U$ 是同构。
由于 $q$ 在对象上形式光滑，存在 $\mathcal{Y}$ 在 $U'$ 上的对象
$y'$、同构 $\beta : y'|_U \to p(x)$ 以及同构
$\gamma' : q(y') \to z'$，使 (\ref{criteria-equation-formally-smooth}) 交换。
由于 $p$ 在对象上形式光滑，存在 $\mathcal{X}$ 在 $U'$ 上的对象
$x'$、同构 $\beta' : x'|_U \to x$ 以及同构
$\gamma'' : p(x') \to y'$，使 (\ref{criteria-equation-formally-smooth}) 交换。
所求解取为 $U'$ 上的 $x'$，并配以同构
$$
q(p(x')) \xrightarrow{q(\gamma'')} q(y') \xrightarrow{\gamma'} z'
$$
和同构 $\beta' : x'|_U \to x$。我们略去
(\ref{criteria-equation-formally-smooth}) 交换性的验证。
\end{proof}

\noindent
注意，代数空间的形式光滑态射类在任意基变换下稳定，并且在目标上相对于
fpqc 拓扑是局部的；参见《空间态射进阶》引理
\ref{spaces-more-morphisms-lemma-base-change-formally-smooth} 和
\ref{spaces-more-morphisms-lemma-descending-property-formally-smooth}。
因此，下一个引理中的条件 (2) 是有意义的。

\begin{lemma}
\label{criteria-lemma-representable-by-spaces-formally-smooth}
设 $p : \mathcal{X} \to \mathcal{Y}$ 是 $(\Sch/S)_{fppf}$ 上
群胚纤维化范畴的 $1$-态射。如果 $p$ 可由代数空间表示，
则下列条件等价：
\begin{enumerate}
\item $p$ 在对象上形式光滑；
\item $p$ 形式光滑（参见《代数叠》定义
\ref{algebraic-definition-relative-representable-property}）。
\end{enumerate}
\end{lemma}

\begin{proof}
假设 (2)。设 $U \subset U'$ 是 $S$ 上仿射概形的一阶加厚，
$y'$ 是 $\mathcal{Y}$ 在 $U'$ 上的对象，$x$ 是 $\mathcal{X}$ 在
$U$ 上的对象，而 $\gamma : p(x) \to y'|_U$ 是同构。
记 $X_{y'}$ 为表示 $2$-纤维积
$$
(\Sch/U')_{fppf} \times_{y', \mathcal{Y}, p} \mathcal{X}.
$$
的 $U'$ 上的代数空间。注意，
$\xi = (U, U \to U', x, \gamma^{-1})$ 定义此 $2$-纤维积在 $U$ 上的
一个对象。由 $2$-Yoneda 引理，$\xi$ 对应于 $U'$ 上的态射
$f_\xi : U \to X_{y'}$。按假设，$X_{y'} \to U'$ 形式光滑，故存在态射
$f' : U' \to X_{y'}$，使 $f_\xi$ 是 $f'$ 与态射 $U \to U'$ 的复合。
$2$-Yoneda 引理还告诉我们，$f'$ 对应于上述 $2$-纤维积在 $U'$ 上的对象
$\xi' = (U', U' \to U', x', \alpha)$，其到 $U$ 的限制恢复 $\xi$。
特别地，我们得到同构
% P11-E0178: frozen source uses |U rather than |_U in the following two displays; preserved canonically.
$\gamma : x'|U \to x$。注意，$\alpha : y' \to p(x')$。
因此，取 $x'$、同构 $\gamma : x'|U \to x$ 以及同构
$\beta = \alpha^{-1} : p(x') \to y'$，即为该问题的一个解。

\medskip\noindent
假设 (1)。选取概形 $T$ 和 $1$-态射
$y : (\Sch/T)_{fppf} \to \mathcal{Y}$。设 $X_y$ 是表示 $2$-纤维积
$(\Sch/T)_{fppf} \times_{y, \mathcal{Y}, p} \mathcal{X}$.
的 $T$ 上的代数空间。我们必须证明 $X_y \to T$ 形式光滑。
因而只需证明：给定 $T$ 上仿射概形的一阶加厚 $U \subset U'$，则
% P11-E0177: frozen source maps X_y(U') to itself here and again below; the formal smoothness criterion normally maps X_y(U') to X_y(U). Preserved canonically.
$X_y(U') \to X_y(U')$ 为满（这里的态射取自 $T$ 上代数空间的范畴）。
令 $y'=y|_{U'}$。由 $2$-Yoneda 引理，$T$ 上的态射
$U \to X_y$ 与二元组 $(x,\alpha)$ 的同构类一一对应，其中 $x$ 是
$\mathcal{X}$ 在 $U$ 上的对象，而 $\alpha : y|_U \to p(x)$ 是同构。
当然，在取逆的意义下，给出 $\alpha$ 等同于给出同构
$\gamma : p(x) \to y'|_U$。$T$ 上的态射 $U' \to X_y$ 情形类似。
因此，(1) 保证 $X_y(U') \to X_y(U')$ 在此情形下为满，结论得证。
\end{proof}







\section{在对象上为满}
\label{criteria-section-formally-surjective}

\noindent
设 $S$ 是概形，$p : \mathcal{X} \to \mathcal{Y}$ 是
$(\Sch/S)_{fppf}$ 上群胚纤维化范畴的 $1$-态射。如果下列条件成立，
则称 $p$ {\it 在对象上为满}：给定由下列各项组成的任意数据
\begin{enumerate}
\item $S$ 上的域 $k$；以及
\item $\mathcal{Y}$ 在 $\Spec(k)$ 上的对象 $y$，
\end{enumerate}
则存在 $S$ 上的域扩张 $K/k$ 以及 $\mathcal{X}$ 在 $\Spec(K)$ 上的对象
$x$，使 $p(x) \cong y|_{\Spec(K)}$。

\begin{lemma}
\label{criteria-lemma-base-change-surjective}
设 $p : \mathcal{X} \to \mathcal{Y}$ 和
$q : \mathcal{Z} \to \mathcal{Y}$ 是 $(\Sch/S)_{fppf}$ 上
群胚纤维化范畴的 $1$-态射。如果
$p : \mathcal{X} \to \mathcal{Y}$ 在对象上为满，则 $p$ 沿 $q$ 的基变换
$p' : \mathcal{X} \times_\mathcal{Y} \mathcal{Z} \to \mathcal{Z}$
也在对象上为满。
\end{lemma}

\begin{proof}
这是形式的。设 $z$ 是 $\mathcal{Z}$ 在域 $k$ 上的对象。
由于 $p$ 在对象上为满，存在扩张 $K/k$、$\mathcal{X}$ 在 $K$ 上的对象
$x$ 以及同构 $\alpha : p(x) \to q(z)|_{\Spec(K)}$。于是
$w = (\Spec(K), x, z|_{\Spec(K)}, \alpha)$ 是
$\mathcal{X} \times_\mathcal{Y} \mathcal{Z}$ 在 $K$ 上的对象，且
$p'(w) = z|_{\Spec(K)}$。
\end{proof}

\begin{lemma}
\label{criteria-lemma-composition-surjective}
设 $p : \mathcal{X} \to \mathcal{Y}$ 和
$q : \mathcal{Y} \to \mathcal{Z}$ 是 $(\Sch/S)_{fppf}$ 上
群胚纤维化范畴的 $1$-态射。如果 $p$ 和 $q$ 均在对象上为满，
则复合 $q \circ p$ 也在对象上为满。
\end{lemma}

\begin{proof}
这是形式的。设 $z$ 是 $\mathcal{Z}$ 在域 $k$ 上的对象。
由于 $q$ 在对象上为满，存在域扩张 $K/k$ 和 $\mathcal{Y}$ 在 $K$ 上的对象
$y$，使
% P11-E0179: frozen source uses x below before x is introduced; z is evidently the intended object, but the canonical formula is preserved.
$q(y) \cong x|_{\Spec(K)}$。由于 $p$ 在对象上为满，存在域扩张 $L/K$
和 $\mathcal{X}$ 在 $L$ 上的对象 $x$，使
$p(x) \cong y|_{\Spec(L)}$。于是域扩张 $L/k$ 和 $\mathcal{X}$ 在
$L$ 上的对象 $x$ 如愿满足 $q(p(x)) \cong z|_{\Spec(L)}$。
\end{proof}

\begin{lemma}
\label{criteria-lemma-representable-by-spaces-surjective}
设 $p : \mathcal{X} \to \mathcal{Y}$ 是 $(\Sch/S)_{fppf}$ 上
群胚纤维化范畴的 $1$-态射。如果 $p$ 可由代数空间表示，
则下列条件等价：
\begin{enumerate}
\item $p$ 在对象上为满；
\item $p$ 为满（参见《代数叠》定义
\ref{algebraic-definition-relative-representable-property}）。
\end{enumerate}
\end{lemma}

\begin{proof}
假设 (2)。设 $k$ 是域，$y$ 是 $\mathcal{Y}$ 在 $k$ 上的对象。
记 $X_y$ 为表示 $2$-纤维积
$$
(\Sch/\Spec(k))_{fppf} \times_{y, \mathcal{Y}, p} \mathcal{X}.
$$
的 $k$ 上的代数空间。由于已假设 $p$ 为满，可知 $X_y$ 非空。
因此可找到域扩张 $K/k$ 和 $X_y$ 的一个 $K$-值点 $x$。
由 $2$-Yoneda 引理，它对应于 $\mathcal{X}$ 在 $K$ 上的对象 $x$ 以及同构
$p(x) \cong y|_{\Spec(K)}$，故 (1) 成立。

\medskip\noindent
假设 (1)。选取概形 $T$ 和 $1$-态射
$y : (\Sch/T)_{fppf} \to \mathcal{Y}$。设 $X_y$ 是表示 $2$-纤维积
$(\Sch/T)_{fppf} \times_{y, \mathcal{Y}, p} \mathcal{X}$.
的 $T$ 上的代数空间。我们必须证明 $X_y \to T$ 为满。
由《空间的态射》定义 \ref{spaces-morphisms-definition-surjective}，
必须证明 $|X_y| \to |T|$ 为满。这恰好意味着：给定 $T$ 上的域 $k$
和态射 $t : \Spec(k) \to T$，存在域扩张 $K/k$ 和态射
$x : \Spec(K) \to X_y$，使得
$$
\xymatrix{
\Spec(K) \ar[d] \ar[r]_x & X_y \ar[d] \\
\Spec(k) \ar[r]^t & T
}
$$
交换。由 $2$-Yoneda 引理，这恰好意味着必须找到 $k \subset K$ 和
$\mathcal{X}$ 在 $K$ 上的对象 $x$，使
$p(x) \cong t^*y|_{\Spec(K)}$。因此 (1) 保证了所需结论。
\end{proof}












\section{代数态射}
\label{criteria-section-algebraic}

\noindent
下列概念有时很有用。

\begin{definition}
\label{criteria-definition-algebraic}
设 $S$ 是概形，$F : \mathcal{X} \to \mathcal{Y}$ 是
$(\Sch/S)_{fppf}$ 上群胚叠的 $1$-态射。如果对每个概形 $T$ 及
$\mathcal{Y}$ 在 $T$ 上的每个对象 $\xi$，$2$-纤维积
$$
(\Sch/T)_{fppf} \times_{\xi, \mathcal{Y}} \mathcal{X}
$$
都是 $S$ 上的代数叠，则称 $F$ 是{\it 代数的}。
\end{definition}

\noindent
采用这一术语后，得到下列推广《代数叠》引理
\ref{algebraic-lemma-representable-morphism-to-algebraic} 的结果。

\begin{lemma}
\label{criteria-lemma-algebraic-morphism-to-algebraic}
设 $S$ 是概形。设 $F : \mathcal{X} \to \mathcal{Y}$ 是
$(\Sch/S)_{fppf}$ 上群胚叠的 $1$-态射。如果
\begin{enumerate}
\item $\mathcal{Y}$ 是代数叠；
\item $F$ 是代数的（见上文），
\end{enumerate}
则 $\mathcal{X}$ 是代数叠。
\end{lemma}

\begin{proof}
由假设 (1)，存在概形 $T$ 和 $\mathcal{Y}$ 在 $T$ 上的对象 $\xi$，
使相应的 $1$-态射 $\xi : (\Sch/T)_{fppf} \to \mathcal{Y}$
% P11-E0180: frozen source says "smooth an surjective"; "and" is evidently intended. Meaning is translated without changing any formula.
光滑且为满。于是
$\mathcal{U} = (\Sch/T)_{fppf} \times_{\xi, \mathcal{Y}} \mathcal{X}$
由假设 (2) 是代数叠。选取概形 $U$ 和满光滑的 $1$-态射
$(\Sch/U)_{fppf} \to \mathcal{U}$。投影
$\mathcal{U} \longrightarrow \mathcal{X}$ 是态射
$\xi : (\Sch/T)_{fppf} \to \mathcal{Y}$ 的基变换，因而满且光滑；
参见《代数叠》引理
\ref{algebraic-lemma-base-change-representable-transformations-property}。
于是复合
$(\Sch/U)_{fppf} \to \mathcal{U} \to \mathcal{X}$
作为满光滑态射的复合，仍然满且光滑；参见《代数叠》引理
\ref{algebraic-lemma-composition-representable-transformations-property}。
因此，由《代数叠》引理
\ref{algebraic-lemma-smooth-surjective-morphism-implies-algebraic}，
$\mathcal{X}$ 是代数叠。
\end{proof}

\begin{lemma}
\label{criteria-lemma-map-from-algebraic}
设 $S$ 是概形，$F : \mathcal{X} \to \mathcal{Y}$ 是
$(\Sch/S)_{fppf}$ 上群胚叠的 $1$-态射。如果 $\mathcal{X}$ 是代数叠，
且 $\Delta : \mathcal{Y} \to \mathcal{Y} \times \mathcal{Y}$
可由代数空间表示，则 $F$ 是代数的。
\end{lemma}

\begin{proof}
选取可表示群胚叠 $\mathcal{U}$ 和满光滑的 $1$-态射
$\mathcal{U} \to \mathcal{X}$。设 $T$ 是概形，$\xi$ 是 $\mathcal{Y}$
在 $T$ 上的对象。$2$-纤维积之间的态射
$$
(\Sch/T)_{fppf} \times_{\xi, \mathcal{Y}} \mathcal{U}
\longrightarrow
(\Sch/T)_{fppf} \times_{\xi, \mathcal{Y}} \mathcal{X}
$$
是 $\mathcal{U} \to \mathcal{X}$ 的基变换，故可由代数空间表示、为满且光滑；
参见《代数叠》引理
\ref{algebraic-lemma-base-change-representable-by-spaces} 和
\ref{algebraic-lemma-base-change-representable-transformations-property}。
由关于 $\mathcal{Y}$ 的对角态射的条件，此态射的源可由代数空间表示；
参见《代数叠》引理 \ref{algebraic-lemma-representable-diagonal}。
因此，由《代数叠》引理
\ref{algebraic-lemma-smooth-surjective-morphism-implies-algebraic}，
其目标是代数叠。
\end{proof}

\begin{lemma}
\label{criteria-lemma-diagonals-and-algebraic-morphisms}
设 $S$ 是概形，$F : \mathcal{X} \to \mathcal{Y}$ 是
$(\Sch/S)_{fppf}$ 上群胚叠的 $1$-态射。如果 $F$ 是代数的，且
$\Delta : \mathcal{Y} \to \mathcal{Y} \times \mathcal{Y}$
可由代数空间表示，则
$\Delta : \mathcal{X} \to \mathcal{X} \times \mathcal{X}$
可由代数空间表示。
\end{lemma}

\begin{proof}
假设 $F$ 是代数的，且
$\Delta : \mathcal{Y} \to \mathcal{Y} \times \mathcal{Y}$
可由代数空间表示。取 $S$ 上的概形 $U$ 以及 $\mathcal{X}$ 在 $U$ 上的
两个对象 $x_1,x_2$。须证明 $\mathit{Isom}(x_1, x_2)$ 是 $U$ 上的代数空间；
参见《代数叠》引理 \ref{algebraic-lemma-representable-diagonal}。
令 $y_i=F(x_i)$。存在集合层的态射
$$
f : \mathit{Isom}(x_1, x_2) \to \mathit{Isom}(y_1, y_2)
$$
，而按假设，其目标是代数空间。因此只需证明 $f$ 可由代数空间表示；
参见《自举》引理
\ref{bootstrap-lemma-representable-by-spaces-over-space}。
所以可以选取 $U$ 上的概形 $V$ 和同构
$\beta : y_{1, V} \to y_{2, V}$，并且须证明函子
$$
(\Sch/V)_{fppf} \to \textit{Sets},\quad
T/V \mapsto \{\alpha : x_{1, T} \to x_{2, T}
\text{ 于 }\mathcal{X}_T \mid F(\alpha) = \beta|_T\}
$$
是代数空间。考虑
$$
\mathcal{Z} = (\Sch/V)_{fppf} \times_{y_{1, V}, \mathcal{Y}} \mathcal{X}
$$
的对象 $z_1 = (V, x_{1, V}, \text{id})$ 和
$z_2 = (V, x_{2, V}, \beta)$。直接验证可知，上述函子在
$(\Sch/V)_{fppf}$ 上等于 $\mathit{Isom}(z_1,z_2)$。
因此，由 $F$ 是代数的这一假设（以及代数叠的定义），它是代数空间。
\end{proof}
















\section{截面空间}
\label{criteria-section-spaces-sections}

\noindent
给定态射 $W \to Z \to U$，可以考虑如下函子：它把 $U$ 上的概形
$U'$ 映为态射 $W \to Z$ 的基变换 $W_{U'} \to Z_{U'}$ 的截面
$\sigma : Z_{U'} \to W_{U'}$ 所成的集合。本节证明关于此函子的若干预备引理。

\begin{lemma}
\label{criteria-lemma-surjection-space-of-sections}
设 $Z \to U$ 是概形的有限态射，$W$ 是代数空间，而
$W \to Z$ 是满平展态射。则存在满平展态射 $U' \to U$ 和截面
$$
\sigma : Z_{U'} \to W_{U'}
$$
，它是态射 $W_{U'} \to Z_{U'}$ 的截面。
\end{lemma}

\begin{proof}
可以选取分离概形 $W'$ 和满平展态射 $W' \to W$。因此，以 $W'$ 替换
$W$ 后，可以假设 $W$ 是分离概形。记 $f : W \to Z$ 且
$\pi : Z \to U$。由于 $W$ 分离，
% P11-E0181: frozen source writes f \circ \pi : W \to U although the typed composition is \pi \circ f; preserved canonically.
$f \circ \pi : W \to U$ 分离（参见《概形》引理
\ref{schemes-lemma-compose-after-separated}）。设 $u \in U$ 是一点。
显然，只需找到 $(U,u)$ 的平展邻域 $(U',u')$，使 $U'$ 上存在截面
$\sigma$。设 $z_1,\ldots,z_r$ 是 $Z$ 中位于 $u$ 上方的点。
对每个 $i$，选取映到 $z_i$ 的点 $w_i\in W$。可以选取平展邻域
$(U', u') \to (U, u)$，使《态射进阶》引理
\ref{more-morphisms-lemma-etale-splits-off-quasi-finite-part-technical-variant}
的结论同时适用于 $Z \to U$ 与诸点 $z_1,\ldots,z_r$，以及
$W \to U$ 与诸点 $w_1,\ldots,w_r$。因此，以 $(U',u')$ 替换
$(U,u)$ 并重新编号后，可以假设所有域扩张
$\kappa(z_i)/\kappa(u)$ 和 $\kappa(w_i)/\kappa(u)$ 都是纯不可分的，
而且存在由开闭子概形给出的不交并分解
$$
Z = V_1 \amalg \ldots \amalg V_r \amalg A, \quad
W = W_1 \amalg \ldots \amalg W_r \amalg B
$$
，其中 $z_i\in V_i$、$w_i\in W_i$，且 $V_i\to U$、$W_i\to U$ 有限。
以 $U \setminus \pi(A)$ 替换 $U$ 后，可以假设 $A=\emptyset$，即
$Z = V_1 \amalg \ldots \amalg V_r$。再以
$W_i \cap f^{-1}(V_i)$ 替换 $W_i$，并以
$B \cup \bigcup W_i \cap f^{-1}(Z \setminus V_i)$ 替换 $B$，
可以假设 $f$ 把 $W_i$ 映入 $V_i$。此时
$f_i = f|_{W_i} : W_i \to V_i$ 是 $U$ 上有限的概形态射，因而有限
（参见《态射》引理 \ref{morphisms-lemma-finite-permanence}）。
它还按假设为平展，满足 $f_i^{-1}(\{z_i\})=w_i$，并诱导剩余域同构
$\kappa(z_i) = \kappa(w_i)$（因为二者都是 $\kappa(u)$ 的纯不可分扩张，
而 $f$ 平展，所以 $\kappa(w_i)/\kappa(z_i)$ 可分）。因此，由
《平展态射》引理 \ref{etale-lemma-finite-etale-one-point}，$f_i$ 在
$z_i$ 的某个邻域 $V_i'$ 上是同构。由于 $\pi : Z \to U$ 是闭映射，
缩小 $U$ 后，可以假设 $W_i \to V_i$ 是同构。这就证明了引理。
\end{proof}

\begin{lemma}
\label{criteria-lemma-space-of-sections}
设 $Z \to U$ 是概形的有限局部自由态射，$W$ 是代数空间，而
$W \to Z$ 是平展态射。则函子
$$
F : (\Sch/U)_{fppf}^{opp} \longrightarrow \textit{Sets},
$$
由规则
$$
U' \longmapsto
F(U') =
\{\sigma : Z_{U'} \to W_{U'}\text{ 下列对象的截面： }W_{U'} \to Z_{U'}\}
$$
定义；它是代数空间，且态射 $F \to U$ 平展。
\end{lemma}

\begin{proof}
先假设 $W \to Z$ 还分离。设 $U'$ 是 $U$ 上的概形且
$\sigma \in F(U')$。由《空间的态射》引理
\ref{spaces-morphisms-lemma-section-immersion}，态射 $\sigma$ 是闭浸入。
此外，由《空间的性质》引理
\ref{spaces-properties-lemma-etale-permanence}，$\sigma$ 平展。
故 $\sigma$ 也是开浸入；参见《空间的态射》引理
\ref{spaces-morphisms-lemma-etale-universally-injective-open}。
换言之，$Z_\sigma = \sigma(Z_{U'}) \subset W_{U'}$ 是开子空间，
使态射 $Z_\sigma \to Z_{U'}$ 为同构。特别地，态射
$Z_\sigma \to U'$ 有限。因此得到函子之间的变换
$$
F \longrightarrow (W/U)_{fin}, \quad
\sigma \longmapsto (U' \to U, Z_\sigma)
$$
，其中 $(W/U)_{fin}$ 是《空间群胚进阶》第
\ref{spaces-more-groupoids-section-finite} 节所引入的态射 $W \to U$
的有限部分。显然，此函子变换为单（因为可从 $Z_\sigma$ 恢复
$\sigma$，即取同构 $Z_\sigma \to Z_{U'}$ 的逆）。由《空间群胚进阶》命题
\ref{spaces-more-groupoids-proposition-finite-algebraic-space}
，$(W/U)_{fin}$ 是 $U$ 上平展的代数空间。因此，为完成此情形的证明，
只需证明 $F \to (W/U)_{fin}$ 可表示且为开浸入。为此，假设给定概形态射
$U' \to U$ 和开子空间 $Z' \subset W_{U'}$，且 $Z' \to U'$ 有限。
只需证明存在开子概形 $U'' \subset U'$，使态射 $T \to U'$ 通过
$U''$ 分解，当且仅当 $Z' \times_{U'} T$ 同构地映到
$Z \times_{U'} T$。这由《空间态射进阶》引理
\ref{spaces-more-morphisms-lemma-where-isomorphism}
推出（这里用到
% P11-E0182: frozen source refers to an undefined B in Z \to B; preserved canonically.
$Z \to B$ 平坦、局部有限表示且有限）。由此已证明 $W \to Z$
既分离又平展的情形。

\medskip\noindent
在一般情形下，选取分离概形 $W'$ 和满平展态射 $W' \to W$。
注意，由于其源分离，态射 $W' \to W$ 和
% P11-E0183: frozen source also claims W \to Z is separated "as their source is separated" although W was not assumed separated in this general case; preserved canonically.
$W \to Z$ 都分离。记 $F'$ 为引理中与 $W' \to Z \to U$ 相伴的函子。
在证明的第一段中，我们已证明 $F'$ 可由 $U$ 上平展的代数空间表示。
由引理 \ref{criteria-lemma-surjection-space-of-sections}，函子映射
$F' \to F$ 对 $\Sch/U$ 上的平展拓扑为满。此外，如果 $U'$ 和
$\sigma : Z_{U'} \to W_{U'}$ 定义一点 $\xi \in F(U')$，则纤维积
$$
F'' = F' \times_{F, \xi} U'
$$
是 $\Sch/U'$ 上与态射
$$
W'_{U'} \times_{W_{U'}, \sigma} Z_{U'} \to Z_{U'} \to U'.
$$
相伴的函子。由于第一个态射是分离态射的基变换，因而分离；由第一段的
结果，$F''$ 是 $U'$ 上平展的代数空间。于是 $F' \to F$ 是满平展的
函子变换，并可由代数空间表示。因此，由《自举》定理
\ref{bootstrap-theorem-final-bootstrap}，$F$ 是代数空间。
由于 $F' \to F$ 是代数空间的满平展态射，而 $F' \to U$ 平展，
所以 $F \to U$ 平展。
\end{proof}









\section{相对态射}
\label{criteria-section-relative-morphisms}

\noindent
继续《态射进阶》第 \ref{more-morphisms-section-relative-morphisms} 节
开始的讨论。

\medskip\noindent
设 $S$ 是概形，$Z \to B$ 和 $X \to B$ 是 $S$ 上代数空间的态射。
给定概形 $T$，可以考虑二元组 $(a,b)$，其中 $a : T \to B$ 是态射，
而 $b : T \times_{a, B} Z \to T \times_{a, B} X$ 是 $T$ 上的态射。
图示为
\begin{equation}
\label{criteria-equation-hom}
\vcenter{
\xymatrix{
T \times_{a, B} Z \ar[rd] \ar[rr]_b & &
T \times_{a, B} X \ar[ld] & Z \ar[rd] & & X \ar[ld] \\
& T \ar[rrr]^a & & & B
}
}
\end{equation}
当然，也可以把 $b$ 看作满足下图交换的态射
$b : T \times_{a, B} Z \to X$：
$$
\xymatrix{
T \times_{a, B} Z \ar[r] \ar[d] \ar@/^1pc/[rrr]_-b &
Z \ar[rd] & & X \ar[ld] \\
T \ar[rr]^a & & B
}
$$
在这种情形下，可以定义函子
\begin{equation}
\label{criteria-equation-hom-functor}
\mathit{Mor}_B(Z, X) : (\Sch/S)^{opp} \longrightarrow \textit{Sets},
\quad
T \longmapsto \{(a, b)\text{ 如上}\}
\end{equation}
有时把它视为定义在 $B$ 上概形范畴上的函子，此时从记号中省去 $a$。

\begin{lemma}
\label{criteria-lemma-hom-functor-sheaf}
设 $S$ 是概形，$Z \to B$ 和 $X \to B$ 是 $S$ 上代数空间的态射。则
\begin{enumerate}
\item $\mathit{Mor}_B(Z, X)$ 是 $(\Sch/S)_{fppf}$ 上的层。
\item 如果 $T$ 是 $S$ 上的代数空间，则存在典范双射
$$
\Mor_{\Sh((\Sch/S)_{fppf})}(T, \mathit{Mor}_B(Z, X))
=
\{(a, b)\text{ 如 }(\ref{criteria-equation-hom})\}
$$
\end{enumerate}
\end{lemma}

\begin{proof}
设 $T$ 是 $S$ 上的代数空间，$\{T_i \to T\}$ 是 $T$ 的一个 fppf 覆盖
（如《空间上的拓扑》第 \ref{spaces-topologies-section-fppf} 节）。
假设 $(a_i, b_i) \in \mathit{Mor}_B(Z, X)(T_i)$，且对所有 $i,j$，
$(a_i, b_i)|_{T_i \times_T T_j} = (a_j, b_j)|_{T_i \times_T T_j}$。
由《空间上的下降》引理
\ref{spaces-descent-lemma-fpqc-universal-effective-epimorphisms}，
存在唯一态射 $a : T \to B$，使 $a_i$ 是 $T_i \to T$ 与 $a$ 的复合。
于是 $\{T_i \times_{a_i, B} Z \to T \times_{a, B} Z\}$ 也是 fppf 覆盖；
同一引理表明存在唯一态射
$b : T \times_{a, B} Z \to T \times_{a, B} X$，使 $b_i$ 是
$T_i \times_{a_i, B} Z \to T \times_{a, B} Z$ 与 $b$ 的复合。
因此，$(a, b) \in \mathit{Mor}_B(Z, X)(T)$ 在 $T_i$ 上限制为
$(a_i,b_i)$，对所有 $i$ 均如此。

\medskip\noindent
注意，上一段的结果特别蕴含 (1)。

\medskip\noindent
设 $T$ 是 $S$ 上的代数空间。为证明 (2)，将在所示集合之间构造互逆映射。
下文所说的“二元组”均指符合 (\ref{criteria-equation-hom}) 的二元组 $(a,b)$。

\medskip\noindent
设 $v : T \to \mathit{Mor}_B(Z, X)$ 是自然变换。选取概形 $U$ 和满平展态射
$p : U \to T$。则 $v(p) \in \mathit{Mor}_B(Z, X)(U)$ 对应于 $U$ 上的
二元组 $(a_U,b_U)$。令 $R = U \times_T U$，投影为 $t,s:R\to U$。
由于 $v$ 是函子变换，$(a_U,b_U)$ 沿 $s$ 和 $t$ 的拉回相同。因此，
由于 $\{U \to T\}$ 是 fppf 覆盖，可以应用第一段的结果，
% P11-E0184: frozen source says "apply ... that deduce"; a coordinating conjunction is missing. Meaning is translated without changing formulas.
从而推得 $T$ 上存在唯一的二元组 $(a,b)$。

\medskip\noindent
反过来，设 $(a,b)$ 是 $T$ 上的二元组。令 $U \to T$、
$R = U \times_T U$ 和 $t,s:R\to U$ 如上。由 Yoneda 引理
（《范畴》引理 \ref{categories-lemma-yoneda}），限制 $(a,b)|_U$
给出函子变换 $v : h_U \to \mathit{Mor}_B(Z, X)$。由于两个拉回
$s^*(a, b)|_U$ 和 $t^*(a, b)|_U$ 相等，可知 $v$ 使两个映射
$h_t, h_s : h_R \to h_U$ 余等。由《空间》引理
\ref{spaces-lemma-space-presentation}，$T=U/R$ 是 fppf 商层；又由 (1)，
$\mathit{Mor}_B(Z, X)$ 是 fppf 层，故 $v$ 通过映射
$T \to \mathit{Mor}_B(Z, X)$ 分解。

\medskip\noindent
我们略去上述两个构造互逆的验证。
\end{proof}

\begin{lemma}
\label{criteria-lemma-base-change-hom-functor}
设 $S$ 是概形，$Z \to B$、$X \to B$ 和 $B' \to B$ 是 $S$ 上
代数空间的态射。令 $Z' = B' \times_B Z$ 且 $X' = B' \times_B X$。则
$$
\mathit{Mor}_{B'}(Z', X')
=
B' \times_B \mathit{Mor}_B(Z, X)
$$
在 $\Sh((\Sch/S)_{fppf})$ 中成立。
\end{lemma}

\begin{proof}
函子层面的等式立即由定义得到。由引理
\ref{criteria-lemma-hom-functor-sheaf}，等式两边都是层；又因为层的纤维积
等于相应预层（即函子）的纤维积，故层层面的等式也随之成立。
\end{proof}

\begin{lemma}
\label{criteria-lemma-etale-covering-hom-functor}
设 $S$ 是概形，$Z \to B$ 和 $X' \to X \to B$ 是 $S$ 上代数空间的态射。
假设
\begin{enumerate}
\item $X' \to X$ 平展；
\item $Z \to B$ 有限局部自由。
\end{enumerate}
则 $\mathit{Mor}_B(Z, X') \to \mathit{Mor}_B(Z, X)$ 可由代数空间表示且平展。
如果 $X' \to X$ 还为满，则
$\mathit{Mor}_B(Z, X') \to \mathit{Mor}_B(Z, X)$ 为满。
\end{lemma}

\begin{proof}
设 $U$ 是概形，且 $\xi=(a,b)$ 是 $\mathit{Mor}_B(Z,X)(U)$ 的元素。
须证明函子
$$
h_U \times_{\xi, \mathit{Mor}_B(Z, X)} \mathit{Mor}_B(Z, X')
$$
可由 $U$ 上平展的代数空间表示。令
$Z_U = U \times_{a, B} Z$ 且 $W = Z_U \times_{b, X} X'$。
则 $W \to Z_U \to U$ 如引理 \ref{criteria-lemma-space-of-sections} 中所述，
而该处定义的层 $F$ 与上述纤维积等同。这证明引理的第一个断言。
第二个断言由此以及引理 \ref{criteria-lemma-surjection-space-of-sections} 推出；
后者保证上述情形中的 $F \to U$ 为满。
\end{proof}

\begin{proposition}
\label{criteria-proposition-hom-functor-algebraic-space}
设 $S$ 是概形，$Z \to B$ 和 $X \to B$ 是 $S$ 上代数空间的态射。
如果 $Z \to B$ 有限局部自由，则 $\mathit{Mor}_B(Z, X)$ 是代数空间。
\end{proposition}

\begin{proof}
选取概形 $B'=\coprod B'_i$，它是仿射概形 $B'_i$ 的不交并，并选取
满平展态射 $B'\to B$。还可以假设 $B'_i\times_B Z$ 是某个环的谱，
该环作为 $\Gamma(B'_i, \mathcal{O}_{B'_i})$-模有限自由。由引理
\ref{criteria-lemma-base-change-hom-functor} 和《空间》引理
\ref{spaces-lemma-base-change-representable-transformations-property}
，态射 $\mathit{Mor}_{B'}(Z', X') \to \mathit{Mor}_B(Z, X)$ 满且平展。
因此，由《自举》定理 \ref{bootstrap-theorem-final-bootstrap}，只需证明
$B=B'$ 是仿射概形 $B'_i$ 的不交并、且每个 $B'_i\times_B Z$ 在
$B'_i$ 上有限自由的情形。事实上，此时只需对每个 $B'_i$ 上的限制证明
该结果。因此，可以假设 $B$ 仿射，且 $\Gamma(Z, \mathcal{O}_Z)$ 是有限自由
$\Gamma(B, \mathcal{O}_B)$-模。

\medskip\noindent
选取由仿射概形不交并而成的概形 $X'$ 以及满平展态射 $X'\to X$。
由引理 \ref{criteria-lemma-etale-covering-hom-functor}，态射
$\mathit{Mor}_B(Z, X') \to \mathit{Mor}_B(Z, X)$ 可由代数空间表示、平展且为满。
因此，由《自举》定理 \ref{bootstrap-theorem-final-bootstrap}，只需证明
$X$ 是仿射概形不交并的情形。这把问题归结为下一段讨论的情形。

\medskip\noindent
假设 $X=\coprod_{i\in I}X_i$ 是仿射概形的不交并，$B$ 仿射，且
$\Gamma(Z, \mathcal{O}_Z)$ 是有限自由 $\Gamma(B, \mathcal{O}_B)$-模。
对任意有限子集 $E\subset I$，令
$$
F_E = \mathit{Mor}_B(Z, \coprod\nolimits_{i \in E} X_i).
$$
由《态射进阶》引理
\ref{more-morphisms-lemma-hom-from-finite-free-into-affine}，$F_E$ 是代数空间。
考虑态射
$$
\coprod\nolimits_{E \subset I\text{ 有限}} F_E
\longrightarrow
\mathit{Mor}_B(Z, X)
$$
每个态射 $F_E\to\mathit{Mor}_B(Z,X)$ 都是开浸入，因为它恰好是参数化
二元组 $(a,b)$ 的轨迹，其中 $b$ 映入 $X$ 的开子概形
$\coprod\nolimits_{i\in E}X_i$。此外，如果 $T$ 拟紧，则对任意二元组
$(a,b)$，$b$ 的像包含于某个有限 $E\subset I$ 所对应的
$\coprod\nolimits_{i\in E}X_i$ 中。因此，所示箭头实际上是开覆盖，
由此结论成立\footnote{这里略去一些集合论论证。确切地，必须证明 $\coprod F_E$
是代数空间。这是因为 $|I|\leq\text{size}(X)$，并且
$\text{size}(F_E)\leq\text{size}(X)$；后者由《态射进阶》引理
\ref{more-morphisms-lemma-hom-from-finite-free-into-affine} 的证明中
对 $F_E$ 的显式描述推出。略去若干细节。}；这里应用了《空间》引理
\ref{spaces-lemma-glueing-algebraic-spaces}。
\end{proof}










\section{标量限制}
\label{criteria-section-restriction-of-scalars}

\noindent
假设 $X \to Z \to B$ 是 $S$ 上代数空间的态射。给定概形 $T$，
可以考虑二元组 $(a,b)$，其中 $a:T\to B$ 是态射，而
$b:T\times_{a,B}Z\to X$ 是 $Z$ 上的态射。图示为
\begin{equation}
\label{criteria-equation-pairs}
\vcenter{
\xymatrix{
& X \ar[d] \\
T \times_{a, B} Z \ar[d] \ar[ru]^b \ar[r] & Z \ar[d] \\
T \ar[r]^a & B
}
}
\end{equation}
在这种情形下，可以定义函子
\begin{equation}
\label{criteria-equation-restriction-of-scalars}
\text{Res}_{Z/B}(X) : (\Sch/S)^{opp} \longrightarrow \textit{Sets},
\quad
T \longmapsto \{(a, b)\text{ 如上}\}
\end{equation}
有时把它视为定义在 $B$ 上概形范畴上的函子，此时从记号中省去 $a$。

\begin{lemma}
\label{criteria-lemma-restriction-of-scalars-sheaf}
设 $S$ 是概形，$X \to Z \to B$ 是 $S$ 上代数空间的态射。则
\begin{enumerate}
\item $\text{Res}_{Z/B}(X)$ 是 $(\Sch/S)_{fppf}$ 上的层。
\item 如果 $T$ 是 $S$ 上的代数空间，则存在典范双射
$$
\Mor_{\Sh((\Sch/S)_{fppf})}(T, \text{Res}_{Z/B}(X))
=
\{(a, b)\text{ 如 }(\ref{criteria-equation-pairs})\}
$$
\end{enumerate}
\end{lemma}

\begin{proof}
设 $T$ 是 $S$ 上的代数空间，$\{T_i \to T\}$ 是 $T$ 的一个 fppf 覆盖
（如《空间上的拓扑》第 \ref{spaces-topologies-section-fppf} 节）。
假设 $(a_i,b_i)\in\text{Res}_{Z/B}(X)(T_i)$，且对所有 $i,j$，
$(a_i, b_i)|_{T_i \times_T T_j} = (a_j, b_j)|_{T_i \times_T T_j}$。
由《空间上的下降》引理
\ref{spaces-descent-lemma-fpqc-universal-effective-epimorphisms}，
存在唯一态射 $a:T\to B$，使 $a_i$ 是 $T_i\to T$ 与 $a$ 的复合。
于是 $\{T_i \times_{a_i, B} Z \to T \times_{a, B} Z\}$ 也是 fppf 覆盖；
同一引理表明存在唯一态射 $b:T\times_{a,B}Z\to X$，使 $b_i$ 是
$T_i \times_{a_i, B} Z \to T \times_{a, B} Z$ 与 $b$ 的复合。
因此，$(a,b)\in\text{Res}_{Z/B}(X)(T)$ 在 $T_i$ 上限制为
$(a_i,b_i)$，对所有 $i$ 均如此。

\medskip\noindent
注意，上一段的结果特别蕴含 (1)。

\medskip\noindent
设 $T$ 是 $S$ 上的代数空间。为证明 (2)，将在所示集合之间构造互逆映射。
下文所说的“二元组”均指符合 (\ref{criteria-equation-pairs}) 的二元组 $(a,b)$。

\medskip\noindent
设 $v:T\to\text{Res}_{Z/B}(X)$ 是自然变换。选取概形 $U$ 和满平展态射
$p:U\to T$。则 $v(p)\in\text{Res}_{Z/B}(X)(U)$ 对应于 $U$ 上的二元组
$(a_U,b_U)$。令 $R=U\times_TU$，投影为 $t,s:R\to U$。
由于 $v$ 是函子变换，$(a_U,b_U)$ 沿 $s$ 和 $t$ 的拉回相同。因此，
由于 $\{U\to T\}$ 是 fppf 覆盖，可以应用第一段的结果，
% P11-E0185: frozen source repeats the "apply ... that deduce" grammatical defect here; translated to the intended grammar.
从而推得 $T$ 上存在唯一的二元组 $(a,b)$。

\medskip\noindent
反过来，设 $(a,b)$ 是 $T$ 上的二元组。令 $U\to T$、
$R=U\times_TU$ 和 $t,s:R\to U$ 如上。由 Yoneda 引理
（《范畴》引理 \ref{categories-lemma-yoneda}），限制 $(a,b)|_U$
给出函子变换 $v:h_U\to\text{Res}_{Z/B}(X)$。由于两个拉回
$s^*(a, b)|_U$ 和 $t^*(a, b)|_U$ 相等，可知 $v$ 使两个映射
$h_t,h_s:h_R\to h_U$ 余等。由《空间》引理
\ref{spaces-lemma-space-presentation}，$T=U/R$ 是 fppf 商层；又由 (1)，
$\text{Res}_{Z/B}(X)$ 是 fppf 层，故 $v$ 通过映射
$T\to\text{Res}_{Z/B}(X)$ 分解。

\medskip\noindent
我们略去上述两个构造互逆的验证。
\end{proof}

\noindent
当然，层 $\text{Res}_{Z/B}(X)$ 自带函子的自然变换
$\text{Res}_{Z/B}(X)\to B$。下文将不加说明地使用它。

\begin{lemma}
\label{criteria-lemma-etale-base-change-restriction-of-scalars}
设 $S$ 是概形，$X\to Z\to B$ 和 $B'\to B$ 是 $S$ 上代数空间的态射。
令 $Z'=B'\times_BZ$ 且 $X'=B'\times_BX$。则
$$
\text{Res}_{Z'/B'}(X')
=
B' \times_B \text{Res}_{Z/B}(X)
$$
在 $\Sh((\Sch/S)_{fppf})$ 中成立。
\end{lemma}

\begin{proof}
函子层面的等式立即由定义得到。由引理
\ref{criteria-lemma-restriction-of-scalars-sheaf}，等式两边都是层；又因为层的
纤维积等于相应预层（即函子）的纤维积，故层层面的等式也随之成立。
\end{proof}

\begin{lemma}
\label{criteria-lemma-etale-covering-restriction-of-scalars}
设 $S$ 是概形，$X'\to X\to Z\to B$ 是 $S$ 上代数空间的态射。假设
\begin{enumerate}
\item $X'\to X$ 平展；
\item $Z\to B$ 有限局部自由。
\end{enumerate}
则 $\text{Res}_{Z/B}(X')\to\text{Res}_{Z/B}(X)$ 可由代数空间表示且平展。
如果 $X'\to X$ 还为满，则
$\text{Res}_{Z/B}(X')\to\text{Res}_{Z/B}(X)$ 为满。
\end{lemma}

\begin{proof}
设 $U$ 是概形，且 $\xi=(a,b)$ 是 $\text{Res}_{Z/B}(X)(U)$ 的元素。
须证明函子
$$
h_U \times_{\xi, \text{Res}_{Z/B}(X)} \text{Res}_{Z/B}(X')
$$
可由 $U$ 上平展的代数空间表示。令
$Z_U=U\times_{a,B}Z$ 且 $W=Z_U\times_{b,X}X'$。
则 $W\to Z_U\to U$ 如引理 \ref{criteria-lemma-space-of-sections} 中所述，
而该处定义的层 $F$ 与上述纤维积等同。这证明引理的第一个断言。
第二个断言由此以及引理 \ref{criteria-lemma-surjection-space-of-sections} 推出；
后者保证上述情形中的 $F\to U$ 为满。
\end{proof}

\noindent
此时，可以利用上述引理证明：只要 $Z\to B$ 有限局部自由，
$\text{Res}_{Z/B}(X)$ 就是代数空间；其证明几乎完全类似于
$\mathit{Mor}_B(Z,X)$ 是代数空间的证明，参见命题
\ref{criteria-proposition-hom-functor-algebraic-space}。不过，我们将从下一个引理
以及 $\mathit{Mor}_B(Z,X)$ 是代数空间这一事实直接推出该结果。

\begin{lemma}
\label{criteria-lemma-fibre-diagram}
设 $S$ 是概形，$X\to Z\to B$ 是 $S$ 上代数空间的态射。下图
$$
\xymatrix{
\mathit{Mor}_B(Z, X) \ar[r] & \mathit{Mor}_B(Z, Z) \\
\text{Res}_{Z/B}(X) \ar[r] \ar[u] & B \ar[u]_{\text{id}_Z}
}
$$
是 $(\Sch/S)_{fppf}$ 上层的笛卡尔图。
\end{lemma}

\begin{proof}
略。提示：这是代数几何函子观点的一道练习。
\end{proof}

\begin{proposition}
\label{criteria-proposition-restriction-of-scalars-algebraic-space}
设 $S$ 是概形，$X\to Z\to B$ 是 $S$ 上代数空间的态射。
如果 $Z\to B$ 有限局部自由，则 $\text{Res}_{Z/B}(X)$ 是代数空间。
\end{proposition}

\begin{proof}
由命题 \ref{criteria-proposition-hom-functor-algebraic-space}，函子
$\mathit{Mor}_B(Z,X)$ 和 $\mathit{Mor}_B(Z,Z)$ 都是代数空间。
因此，结论由引理 \ref{criteria-lemma-fibre-diagram} 的笛卡尔图以及下述事实推出：
代数空间的纤维积存在，并由其底层集合层范畴中的纤维积给出
（参见《空间》引理
\ref{spaces-lemma-fibre-product-spaces-over-sheaf-with-representable-diagonal}）。
\end{proof}






\section{有限 Hilbert 叠}
\label{criteria-section-finite-hilbert-stacks}

\noindent
本节证明若干关于《叠的例子》第
\ref{examples-stacks-section-hilbert-d-stack} 节所引入的有限 Hilbert 叠
$\mathcal{H}_d(\mathcal{X}/\mathcal{Y})$ 的结果。

\begin{lemma}
\label{criteria-lemma-map-hilbert}
考虑 $2$-交换图
$$
\xymatrix{
\mathcal{X}' \ar[r]_G \ar[d]_{F'} & \mathcal{X} \ar[d]^F \\
\mathcal{Y}' \ar[r]^H & \mathcal{Y}
}
$$
，其中各项为 $(\Sch/S)_{fppf}$ 上的群胚叠，并给定 $2$-同构
$\gamma : H \circ F' \to F \circ G$。在这种情形下，得到典范 $1$-态射
$\mathcal{H}_d(\mathcal{X}'/\mathcal{Y}') \to
\mathcal{H}_d(\mathcal{X}/\mathcal{Y})$.
此态射与《叠的例子》公式
(\ref{examples-stacks-equation-diagram-hilbert-d-stack}) 中的遗忘 $1$-态射相容。
\end{lemma}

\begin{proof}
把对象 $(U, Z, y', x', \alpha')$ 映为对象
$(U, Z, H(y'), G(x'), \gamma \star \text{id}_H \star \alpha')$
，其中 $\star$ 表示 $2$-态射的水平复合；参见《范畴》定义
\ref{categories-definition-horizontal-composition}。对态射
$(f, g, b, a) :
(U_1, Z_1, y_1', x_1', \alpha_1') \to (U_2, Z_2, y_2', x_2', \alpha_2')$
，赋予态射
$(f, g, H(b), G(a))$.
略去它定义 $(\Sch/S)_{fppf}$ 上范畴之间函子的验证。
\end{proof}

\begin{lemma}
\label{criteria-lemma-cartesian-map-hilbert}
在引理 \ref{criteria-lemma-map-hilbert} 的情形下，假设给定方形是 $2$-笛卡尔的。
则图
$$
\xymatrix{
\mathcal{H}_d(\mathcal{X}'/\mathcal{Y}') \ar[r] \ar[d] &
\mathcal{H}_d(\mathcal{X}/\mathcal{Y}) \ar[d] \\
\mathcal{Y}' \ar[r] &
\mathcal{Y}
}
$$
是 $2$-笛卡尔的。
\end{lemma}

\begin{proof}
由引理 \ref{criteria-lemma-map-hilbert} 得到 $2$-交换图，因而得到 $1$-态射
（即函子）
$$
\mathcal{H}_d(\mathcal{X}'/\mathcal{Y}')
\longrightarrow
\mathcal{Y}' \times_\mathcal{Y} \mathcal{H}_d(\mathcal{X}/\mathcal{Y})
$$
说明此函子为何本质满。右侧范畴的对象由下列数据给出：$S$ 上的概形
$U$、$\mathcal{Y}'_U$ 的对象 $y'$、
$\mathcal{H}_d(\mathcal{X}/\mathcal{Y})$ 在 $U$ 上的对象
$(U,Z,y,x,\alpha)$，以及 $\mathcal{Y}_U$ 中的同构 $H(y')\to y$。
假设恰好意味着存在 $\mathcal{X}'_Z$ 的对象 $x'$，并存在与 $\alpha$
相容的同构 $G(x')\cong x$ 和 $\alpha':y'|_Z\to F'(x')$。
于是 $(U,Z,y',x',\alpha')$ 是
$\mathcal{H}_d(\mathcal{X}'/\mathcal{Y}')$ 在 $U$ 上的对象。略去细节。
\end{proof}

\begin{lemma}
\label{criteria-lemma-etale-covering-hilbert}
在引理 \ref{criteria-lemma-map-hilbert} 的情形下，假设
\begin{enumerate}
\item $\mathcal{Y}' = \mathcal{Y}$ 且 $H = \text{id}_\mathcal{Y}$；
\item $G$ 可由代数空间表示且平展。
\end{enumerate}
则 $\mathcal{H}_d(\mathcal{X}'/\mathcal{Y}) \to
\mathcal{H}_d(\mathcal{X}/\mathcal{Y})$ 可由代数空间表示且平展。
如果 $G$ 还为满，则
$\mathcal{H}_d(\mathcal{X}'/\mathcal{Y}) \to
\mathcal{H}_d(\mathcal{X}/\mathcal{Y})$ 为满。
\end{lemma}

\begin{proof}
设 $U$ 是概形，而 $\xi=(U,Z,y,x,\alpha)$ 是
$\mathcal{H}_d(\mathcal{X}/\mathcal{Y})$ 在 $U$ 上的对象。
须证明 $2$-纤维积
\begin{equation}
\label{criteria-equation-to-show}
(\Sch/U)_{fppf}
\times_{\xi, \mathcal{H}_d(\mathcal{X}/\mathcal{Y})}
\mathcal{H}_d(\mathcal{X}'/\mathcal{Y})
\end{equation}
可由 $U$ 上平展的代数空间表示。它在 $U'$ 上的对象对应于
$\mathcal{X}'$ 在 $Z_{U'}$ 上的纤维范畴中的对象 $x'$，并满足
$G(x')\cong x|_{Z_{U'}}$。按假设，$2$-纤维积
$$
(\Sch/Z)_{fppf} \times_{x, \mathcal{X}} \mathcal{X}'
$$
可由代数空间 $W$ 表示，且投影 $W\to Z$ 平展。于是
(\ref{criteria-equation-to-show}) 可由引理 \ref{criteria-lemma-space-of-sections} 中引入的
代数空间 $F$ 表示；该空间参数化 $U$ 上 $W\to Z$ 的截面。
由于 $F\to U$ 平展，可知
$\mathcal{H}_d(\mathcal{X}'/\mathcal{Y}) \to
\mathcal{H}_d(\mathcal{X}/\mathcal{Y})$ 可由代数空间表示且平展。
最后，如果 $\mathcal{X}'\to\mathcal{X}$ 还为满，则 $W\to Z$ 为满，
因而由引理 \ref{criteria-lemma-surjection-space-of-sections}，$F\to U$ 为满。
所以在此情形下
$\mathcal{H}_d(\mathcal{X}'/\mathcal{Y}) \to
\mathcal{H}_d(\mathcal{X}/\mathcal{Y})$ 也为满。
\end{proof}

\begin{lemma}
\label{criteria-lemma-etale-map-hilbert}
在引理 \ref{criteria-lemma-map-hilbert} 的情形下，假设 $G$、$H$ 均可由代数空间
表示且平展。则 $\mathcal{H}_d(\mathcal{X}'/\mathcal{Y}') \to
\mathcal{H}_d(\mathcal{X}/\mathcal{Y})$ 可由代数空间表示且平展。
如果 $H$ 还为满，且诱导函子
$\mathcal{X}' \to \mathcal{Y}' \times_\mathcal{Y} \mathcal{X}$
为满，则
$\mathcal{H}_d(\mathcal{X}'/\mathcal{Y}') \to
\mathcal{H}_d(\mathcal{X}/\mathcal{Y})$ 为满。
\end{lemma}

\begin{proof}
令 $\mathcal{X}'' = \mathcal{Y}' \times_\mathcal{Y} \mathcal{X}$。
由引理 \ref{criteria-lemma-etale-permanence}，$1$-态射
$\mathcal{X}' \to \mathcal{X}''$ 可由代数空间表示且平展（特别地，
引理第二个断言中要求 $\mathcal{X}'\to\mathcal{X}''$ 为满的条件是有意义的）。
得到 $2$-交换图
$$
\xymatrix{
\mathcal{X}' \ar[r] \ar[d] &
\mathcal{X}'' \ar[r] \ar[d] &
\mathcal{X} \ar[d] \\
\mathcal{Y}' \ar[r] &
\mathcal{Y}' \ar[r] &
\mathcal{Y}
}
$$
由引理 \ref{criteria-lemma-cartesian-map-hilbert}，
$\mathcal{H}_d(\mathcal{X}''/\mathcal{Y}')$ 是
$\mathcal{H}_d(\mathcal{X}/\mathcal{Y})$ 沿
$\mathcal{Y}'\to\mathcal{Y}$ 的基变换。特别地，
$\mathcal{H}_d(\mathcal{X}''/\mathcal{Y}') \to
\mathcal{H}_d(\mathcal{X}/\mathcal{Y})$ 可由代数空间表示且平展；参见
《代数叠》引理
\ref{algebraic-lemma-base-change-representable-transformations-property}.
此外，如果 $H$ 为满，它也为满。因此，只要证明该结果对图中的左方形成立，
证明即告完成。这样便归结到 $\mathcal{Y}'=\mathcal{Y}$ 的情形，
而这正是引理 \ref{criteria-lemma-etale-covering-hilbert} 的内容。
\end{proof}

\begin{lemma}
\label{criteria-lemma-relative-hilbert}
设 $F : \mathcal{X} \to \mathcal{Y}$ 是 $(\Sch/S)_{fppf}$ 上群胚叠的
$1$-态射。假设
$\Delta : \mathcal{Y} \to \mathcal{Y} \times \mathcal{Y}$
可由代数空间表示。则态射
$$
\mathcal{H}_d(\mathcal{X}/\mathcal{Y})
\longrightarrow
\mathcal{H}_d(\mathcal{X}) \times \mathcal{Y}
$$
（见《叠的例子》公式
(\ref{examples-stacks-equation-diagram-hilbert-d-stack})）可由代数空间表示。
\end{lemma}

\begin{proof}
设 $U$ 是概形，而 $\xi=(U,Z,p,x,1)$ 是
$\mathcal{H}_d(\mathcal{X})=\mathcal{H}_d(\mathcal{X}/S)$ 在 $U$ 上的对象。
这里 $p$ 只是 $U$ 的结构态射。由于一切都在 $S$ 上，第五个分量 $1$
存在且唯一。另设 $y$ 是 $\mathcal{Y}$ 在 $U$ 上的对象。须证明
$2$-纤维积
\begin{equation}
\label{criteria-equation-res-isom}
(\Sch/U)_{fppf}
\times_{\xi \times y, \mathcal{H}_d(\mathcal{X}) \times \mathcal{Y}}
\mathcal{H}_d(\mathcal{X}/\mathcal{Y})
\end{equation}
可由代数空间表示。为说明这一点，引入
$$
I = \mathit{Isom}_\mathcal{Y}(y|_Z, F(x))
$$
。按假设，它是 $Z$ 上的代数空间。设 $a:U'\to U$ 是 $U$ 上的概形。
给出 (\ref{criteria-equation-res-isom}) 的纤维范畴在 $U'$ 上的对象意味着什么？
这意味着给出 $\mathcal{H}_d(\mathcal{X}/\mathcal{Y})$ 在 $U'$ 上的对象
$\xi'=(U',Z',y',x',\alpha')$，以及同构
$(U', Z', p', x', 1) \cong (U, Z, p, x, 1)|_{U'}$ 和
$y' \cong y|_{U'}$。因此，$\xi'$ 同构于
$(U', U' \times_{a, U} Z, a^*y, x|_{U' \times_{a, U} Z}, \alpha)$
，其中有某个态射
$$
\alpha :
a^*y|_{U' \times_{a, U} Z}
\longrightarrow
F(x|_{U' \times_{a, U} Z})
$$
，它位于 $\mathcal{Y}$ 在 $U'\times_{a,U}Z$ 上的纤维范畴中。
因此，可以把 $\alpha$ 看作态射 $b:U'\times_{a,U}Z\to I$。
由此可知，(\ref{criteria-equation-res-isom}) 可由 $\text{Res}_{Z/U}(I)$ 表示；
由命题 \ref{criteria-proposition-restriction-of-scalars-algebraic-space}，后者是代数空间。
\end{proof}

\noindent
下一个引理是引理 \ref{criteria-lemma-etale-covering-hilbert} 的一个（部分）推广。

\begin{lemma}
\label{criteria-lemma-representable-on-top}
设 $F : \mathcal{X} \to \mathcal{Y}$ 和
$G : \mathcal{X}' \to \mathcal{X}$ 是 $(\Sch/S)_{fppf}$ 上群胚叠的
$1$-态射。如果 $G$ 可由代数空间表示，则 $1$-态射
$$
\mathcal{H}_d(\mathcal{X}'/\mathcal{Y})
\longrightarrow
\mathcal{H}_d(\mathcal{X}/\mathcal{Y})
$$
可由代数空间表示。
\end{lemma}

\begin{proof}
设 $U$ 是概形，而 $\xi=(U,Z,y,x,\alpha)$ 是
$\mathcal{H}_d(\mathcal{X}/\mathcal{Y})$ 在 $U$ 上的对象。
须证明 $2$-纤维积
\begin{equation}
\label{criteria-equation-to-show-again}
(\Sch/U)_{fppf}
\times_{\xi, \mathcal{H}_d(\mathcal{X}/\mathcal{Y})}
\mathcal{H}_d(\mathcal{X}'/\mathcal{Y})
\end{equation}
可由 $U$ 上的代数空间表示。它在 $a:U'\to U$ 上的对象对应于
$\mathcal{X}'$ 在 $U'\times_{a,U}Z$ 上的对象 $x'$，并满足
$G(x')\cong x|_{U'\times_{a,U}Z}$。按假设，$2$-纤维积
$$
(\Sch/Z)_{fppf} \times_{x, \mathcal{X}} \mathcal{X}'
$$
可由 $Z$ 上的代数空间 $W$ 表示。因此，(\ref{criteria-equation-to-show-again})
可由 $\text{Res}_{Z/U}(W)$ 表示；由命题
\ref{criteria-proposition-restriction-of-scalars-algebraic-space}，后者是代数空间。
\end{proof}

\begin{lemma}
\label{criteria-lemma-limit-preserving}
设 $F : \mathcal{X} \to \mathcal{Y}$ 是 $(\Sch/S)_{fppf}$ 上群胚叠的
$1$-态射。假设 $F$ 可由代数空间表示且局部有限表示。则
$$
p : \mathcal{H}_d(\mathcal{X}/\mathcal{Y}) \to \mathcal{Y}
$$
在对象上保极限。
\end{lemma}

\begin{proof}
这意味着必须证明如下命题：给定
\begin{enumerate}
\item 仿射概形 $U = \lim_i U_i$，它写成 $S$ 上仿射概形 $U_i$ 的有向极限；
\item 对某个 $i$，$\mathcal{Y}$ 在 $U_i$ 上的对象 $y_i$；以及
\item $\mathcal{H}_d(\mathcal{X}/\mathcal{Y})$ 在 $U$ 上的对象
$\Xi = (U, Z, y, x, \alpha)$，满足 $y = y_i|_U$，
\end{enumerate}
则存在 $i' \geq i$ 和
$\Xi_{i'} = (U_{i'}, Z_{i'}, y_{i'}, x_{i'}, \alpha_{i'})$，它是
$\mathcal{H}_d(\mathcal{X}/\mathcal{Y})$ 在 $U_{i'}$ 上的对象，并满足
$\Xi_{i'}|_U = \Xi$ 和 $y_{i'} = y_i|_{U_{i'}}$。最后这两个等式恰好保证
(\ref{criteria-equation-limit-preserving}) 交换。

\medskip\noindent
设 $X_{y_i} \to U_i$ 是表示 $2$-纤维积
$$
(\Sch/U_i)_{fppf} \times_{y_i, \mathcal{Y}, F} \mathcal{X}.
$$
的代数空间。由关于 $F$ 的假设，$X_{y_i} \to U_i$ 局部有限表示。
% P11-E0186: frozen source contains the isolated incomplete sentence "Write Xi." here; retained semantically without inventing missing content.
写下 $\Xi$。显然，$\xi = (Z, Z \to U_i, x, \alpha)$ 是上述
$2$-纤维积的对象，故 $\xi$ 给出 $U_i$ 上代数空间的态射
$f_\xi : Z \to X_{y_i}$（因为 $X_{y_i}$ 是
% P11-E0187: frozen source says "isomorphisms classes"; the intended singular attributive form is translated.
$(\Sch/U_i)_{fppf} \times_{y, \mathcal{Y}, F} \mathcal{X}$
的对象同构类函子；参见《代数叠》引理
\ref{algebraic-lemma-characterize-representable-by-space}）。
% P11-E0188: frozen source uses y rather than y_i in the displayed fibre product above; the canonical math token is preserved.
由《极限》引理 \ref{limits-lemma-descend-finite-presentation} 和
\ref{limits-lemma-descend-finite-locally-free}
，存在 $i' \geq i$ 和次数为 $d$ 的有限局部自由态射
$Z_{i'} \to U_{i'}$，其到 $U$ 的基变换为 $Z$。由《空间的极限》命题
\ref{spaces-limits-proposition-characterize-locally-finite-presentation}
，以更大的指标替换 $i'$ 后，可以假设存在态射
$f_{i'} : Z_{i'} \to X_{y_i}$，使下图交换：
$$
\xymatrix{
Z \ar[d] \ar[r] \ar@/^3ex/[rr]^{f_\xi} &
Z_{i'} \ar[d] \ar[r]_{f_{i'}} & X_{y_i} \ar[d] \\
U \ar[r] & U_{i'} \ar[r] & U_i
}
$$
令
$\Xi_{i'} = (U_{i'}, Z_{i'}, y_{i'}, x_{i'}, \alpha_{i'})$，其中
\begin{enumerate}
\item $y_{i'}$ 是 $\mathcal{Y}$ 在 $U_{i'}$ 上的对象，即 $y_i$ 到
$U_{i'}$ 的拉回；
\item $x_{i'}$ 是 $\mathcal{X}$ 在 $Z_{i'}$ 上的对象，它经由
$2$-Yoneda 引理对应于 $1$-态射
$$
(\Sch/Z_{i'})_{fppf} \to
\mathcal{S}_{X_{y_i}} \to
(\Sch/U_i)_{fppf} \times_{y_i, \mathcal{Y}, F} \mathcal{X} \to
\mathcal{X}
$$
，其中中间箭头是定义 $X_{y_i}$ 的等价（记号如《代数叠》第
\ref{algebraic-section-representable-by-algebraic-spaces} 节和第
\ref{algebraic-section-split} 节）。
\item $\alpha_{i'} : y_{i'}|_{Z_{i'}} \to F(x_{i'})$ 是由下图的
$2$-交换性得到的同构：
$$
\xymatrix{
(\Sch/Z_{i'})_{fppf} \ar[r] \ar[rd] &
(\Sch/U_i)_{fppf} \times_{y_i, \mathcal{Y}, F} \mathcal{X}
\ar[r] \ar[d] &
\mathcal{X} \ar[d]^F \\
& (\Sch/U_{i'})_{fppf} \ar[r] & \mathcal{Y}
}
$$
\end{enumerate}
回顾 $f_\xi : Z \to X_{y_i}$ 是与对象
$\xi = (Z, Z \to U_i, x, \alpha)$ 相对应的态射；该对象是
$(\Sch/U_i)_{fppf} \times_{y_i, \mathcal{Y}, F} \mathcal{X}$
在 $Z$ 上的对象。按构造，$f_{i'}$ 是与对象
$\xi_{i'} = (Z_{i'}, Z_{i'} \to U_i, x_{i'}, \alpha_{i'})$ 相对应的态射。
由于 $f_\xi = f_{i'} \circ (Z \to Z_{i'})$，可知对象
$\xi_{i'} = (Z_{i'}, Z_{i'} \to U_i, x_{i'}, \alpha_{i'})$ 到 $Z$ 的拉回为
$\xi$。因此，$x_{i'}$ 拉回为 $x$，而 $\alpha_{i'}$ 拉回为 $\alpha$。
这意味着 $\Xi_{i'}$ 到 $U$ 的拉回为 $\Xi$，结论得证。
\end{proof}










\section{一点的有限 Hilbert 叠}
\label{criteria-section-hilbert-point}

\noindent
设 $d\geq1$ 是整数。在《叠的例子》定义
\ref{examples-stacks-definition-hilbert-d-stack} 中，我们定义了群胚叠
$\mathcal{H}_d$。本节证明 $\mathcal{H}_d$ 是代数叠。全文假设
$S=\Spec(\mathbf{Z})$；一般情形由基变换随之得到。回顾
$\mathcal{H}_d$ 在概形 $T$ 上的纤维范畴，是次数为 $d$ 的有限局部自由态射
$\pi:Z\to T$ 所成的范畴。我们不直接分类这些态射，而先研究拟凝聚代数层
$\pi_*\mathcal{O}_Z$。

\medskip\noindent
设 $R$ 是环。暂作如下定义：一个{\it $R$ 上的自由 $d$ 维代数}
由 $R^{\oplus d}$ 上的交换 $R$-代数结构 $m$ 给出，并要求
$e_1=(1,0,\ldots,0)$ 为单位元\footnote{也许更适合把它看作由乘法映射
$m:R^{\oplus d}\otimes_RR^{\oplus d}\to R^{\oplus d}$ 和满足若干公理的
环映射 $\psi:R\to R^{\oplus d}$ 组成的二元组。}。把 $m$ 看作 $R$-线性映射
$$
m : R^{\oplus d} \otimes_R R^{\oplus d} \longrightarrow R^{\oplus d}
$$
，它满足 $m(e_1,x)=m(x,e_1)=x$，并且 $m$ 定义交换结合的环结构。
若写 $m(e_i,e_j)=\sum a_{ij}^ke_k$，则这归结为条件
$$
\left\{
\begin{matrix}
\sum_l a_{ij}^la_{lk}^m = \sum_l a_{il}^ma_{jk}^l & \forall i, j, k, m \\
a_{ij}^k = a_{ji}^k & \forall i, j, k \\
a_{i1}^j = \delta_{ij} & \forall i, j
\end{matrix}
\right.
$$
，其中 $\delta_{ij}$ 是 Kronecker $\delta$-函数。于是定义
$$
R_{univ} = \mathbf{Z}[a_{ij}^k]/J
$$
，其中理想 $J$ 由上述关系生成。记
$$
m_{univ} :
R_{univ}^{\oplus d} \otimes_{R_{univ}} R_{univ}^{\oplus d}
\longrightarrow
R_{univ}^{\oplus d}
$$
% P11-E0189: frozen source's naming construction "Denote [display] the ..." omits "by"; translated naturally.
为 $R_{univ}$ 上的自由 $d$ 维代数 $m$，其结构常数为 $a_{ij}^k$ 模 $J$
的类。显然，给定环 $R$ 上任意自由 $d$ 维代数 $m$，存在唯一的
$\mathbf{Z}$-代数同态 $\psi:R_{univ}\to R$，使
$\psi_*m_{univ}=m$（这意味着 $m$ 由基变换函子
$-\otimes_{R_{univ}}R$ 作用于 $m_{univ}$ 得到）。换言之，令
$X=\Spec(R_{univ})$，得到典范等同
$$
X(T) = \{\text{自由 }d\text{-维代数 }m\text{ 相对于 }R\}
$$
，其中 $T=\Spec(R)$ 变化。通过 Zariski 局部化，得到下列看似更一般的等同
\begin{equation}
\label{criteria-equation-objects}
X(T) = \{\text{自由 }d\text{-维代数 }
m\text{ 相对于 }\Gamma(T, \mathcal{O}_T)\}
\end{equation}
，对任意概形 $T$ 成立。

\medskip\noindent
接下来略述{\it 自由 $d$ 维 $R$-代数的同构}。设 $m$、$m'$ 是环 $R$
上的两个自由 $d$ 维代数。一个{\it 从 $m$ 到 $m'$ 的同构}由可逆
$R$-线性映射
$$
\varphi : R^{\oplus d} \longrightarrow R^{\oplus d}
$$
给出，并满足 $\varphi(e_1)=e_1$ 以及
$$
m \circ \varphi \otimes \varphi = \varphi \circ m'.
$$
注意，可以复合这些同构，从而使 $R$ 上自由 $d$ 维代数的全体成为范畴。
由此得到函子
\begin{equation}
\label{criteria-equation-FAd}
FA_d : \Sch_{fppf}^{opp} \longrightarrow \textit{Groupoids}
\end{equation}
，它从概形范畴取值于群胚：对概形 $T$，赋予
$\Gamma(T,\mathcal{O}_T)$ 上自由 $d$ 维代数的集合，并用刚才定义的
同构概念赋予其范畴结构。

\medskip\noindent
上述讨论提示我们考虑群值函子 $G$，它把任意概形 $T$ 映为群
$$
G(T) = \{g \in \text{GL}_d(\Gamma(T, \mathcal{O}_T)) \mid g(e_1) = e_1\}
$$
显然，$G\subset\text{GL}_d$（参见《群胚》例
\ref{groupoids-example-general-linear-group}）是由方程
$x_{11}=1$ 和 $x_{i1}=0$（$i>1$）截出的闭子群概形。因此，$G$ 是
$\Spec(\mathbf{Z})$ 上的光滑仿射群概形。考虑作用
$$
a : G \times_{\Spec(\mathbf{Z})} X \longrightarrow X
$$
，它把左侧的 $T$-值点 $(g,m)$（其中 $T=\Spec(R)$）映为如下给出的
$R$ 上自由 $d$ 维代数：
$$
a(g, m) = g^{-1} \circ m \circ g \otimes g.
$$
注意，这意味着 $g$ 定义自由 $d$ 维 $R$-代数的同构
$m\to a(g,m)$。略去 $a$ 确实定义群概形 $G$ 在概形 $X$ 上作用的验证。

\begin{lemma}
\label{criteria-lemma-represent-FAd}
在 (\ref{criteria-equation-FAd}) 中定义的群胚值函子 $FA_d$ 同构（！）于如下
群胚值函子：它把概形 $T$ 映为具有下列数据的范畴
\begin{enumerate}
\item 对象集为 $X(T)$；
\item 态射集为 $G(T) \times X(T)$；
\item $s : G(T) \times X(T) \to X(T)$ 是投影映射；
\item $t : G(T) \times X(T) \to X(T)$ 是 $a(T)$；
\item 复合\newline
$G(T) \times X(T) \times_{s, X(T), t} G(T) \times X(T)
\to G(T) \times X(T)$ 由
$((g, m), (g', m')) \mapsto (gg', m')$ 给出。
\end{enumerate}
\end{lemma}

\begin{proof}
在 (\ref{criteria-equation-objects}) 中已经看到对象上的规则。上文还看到，
对任意自由 $d$ 维代数 $m$，$g\in G(T)$ 可视为从 $m$ 到 $a(g,m)$
的态射。反过来，任意态射 $m\to m'$ 由可逆线性映射 $\varphi$ 给出；
它对应于元素 $g\in G(T)$，并满足 $m'=a(g,m)$。
\end{proof}

\noindent
事实上，上一个引理所述的群胚 $(X,G\times X,s,t,c)$ 正是《群胚》引理
\ref{groupoids-lemma-groupoid-from-action} 中与作用
$a:G\times X\to X$ 相伴的群胚。由于 $G$ 在 $\Spec(\mathbf{Z})$ 上光滑，
可知两个态射 $s,t:G\times X\to X$ 都光滑：由对称性，只需证明其中一个，
而 $s$ 是 $G\to\Spec(\mathbf{Z})$ 的基变换。因此
$(G\times X,X,s,t,c)$ 是光滑群胚概形，并且由《代数叠》定理
\ref{algebraic-theorem-smooth-groupoid-gives-algebraic-stack}，商叠
$[X/G]$ 是代数叠。

\begin{proposition}
\label{criteria-proposition-finite-hilbert-point}
叠 $\mathcal{H}_d$ 等价于上述商叠 $[X/G]$。特别地，$\mathcal{H}_d$
是代数叠。
\end{proposition}

\begin{proof}
注意，由《空间中的群胚》定义
\ref{spaces-groupoids-definition-quotient-stack}，商叠 $[X/G]$ 是与下述
“群胚预层”相伴的群胚纤维化范畴的叠化；该预层把概形 $T$ 映为群胚
$$
(X(T), G(T) \times X(T), s, t, c).
$$
由引理 \ref{criteria-lemma-represent-FAd}，此“群胚预层”同构于 $FA_d$。
因此，只需证明
% P11-E0190: frozen source has the superfluous article "the" before mathcal H_d; translated naturally.
$\mathcal{H}_d$ 是（与“群胚预层”相伴的群胚纤维化范畴）$FA_d$ 的叠化。
为此，先定义函子
$$
\Spec : FA_d \longrightarrow \mathcal{H}_d
$$
回顾 $\mathcal{H}_d$ 在概形 $T$ 上的纤维范畴，是次数为 $d$ 的有限局部自由
态射 $Z\to T$ 所成的范畴。因此，给定概形 $T$ 和自由 $d$ 维
$\Gamma(T,\mathcal{O}_T)$-代数 $m$，可以赋予它对象
$$
Z = \underline{\Spec}_T(\mathcal{A})
$$
，它属于 $\mathcal{H}_{d,T}$；其中
$\mathcal{A}=\mathcal{O}_T^{\oplus d}$ 经由 $m$ 赋有
$\mathcal{O}_T$-代数结构。此外，若 $m'$ 是另一个这样的自由 $d$ 维
$\Gamma(T,\mathcal{O}_T)$-代数，而 $\varphi:m\to m'$ 是二者的同构，
则诱导的 $\mathcal{O}_T$-线性映射
$\varphi : \mathcal{O}_T^{\oplus d} \to \mathcal{O}_T^{\oplus d}$
诱导拟凝聚 $\mathcal{O}_T$-代数的同构
$$
\varphi : \mathcal{A}' \longrightarrow \mathcal{A}
$$
。因此
$$
\underline{\Spec}_T(\varphi) :
\underline{\Spec}_T(\mathcal{A})
\longrightarrow
\underline{\Spec}_T(\mathcal{A}')
$$
是纤维范畴 $\mathcal{H}_{d,T}$ 中的态射。略去此构造与基变换相容的验证，
于是确实得到上述函子 $\Spec:FA_d\to\mathcal{H}_d$。

\medskip\noindent
为证明 $\Spec:FA_d\to\mathcal{H}_d$ 在 $FA_d$ 的叠化与
$\mathcal{H}_d$ 之间诱导等价，只需验证
\begin{enumerate}
\item $\mathit{Isom}(m, m') = \mathit{Isom}(\Spec(m), \Spec(m'))$
，对任意 $m,m'\in FA_d(T)$；
\item 对任意概形 $T$ 和 $\mathcal{H}_{d,T}$ 的任意对象 $Z\to T$，
存在覆盖 $\{T_i\to T\}$，使对某个 $m\in FA_d(T_i)$，$Z|_{T_i}$
同构于 $\Spec(m)$；
\end{enumerate}
参见《叠》引理 \ref{stacks-lemma-stackify-groupoids}。第一个断言由下述观察推出：
任意同构
$$
\underline{\Spec}_T(\mathcal{A})
\longrightarrow
\underline{\Spec}_T(\mathcal{A}')
$$
在 $\mathcal{A}=\mathcal{A}'=\mathcal{O}_T^{\oplus d}$ 作为模时，必由
整体可逆矩阵 $g$ 给出。为证明第二个断言，设 $\pi:Z\to T$ 是次数为
$d$ 的有限局部自由态射。于是，
% P11-E0191: frozen source says "a locally free sheaf O_T-modules"; "of" is missing. Translated naturally.
$\mathcal{A}$ 是秩为 $d$ 的局部自由 $\mathcal{O}_T$-模层。考虑元素
$1\in\Gamma(T,\mathcal{A})$。对每个 $t\in T$，此元素在
$\mathcal{A}\otimes_{\mathcal{O}_{T,t}}\kappa(t)$ 中非零，因为概形
$Z_t = \Spec(\mathcal{A} \otimes_{\mathcal{O}_{T, t}} \kappa(t))$
在 $\kappa(t)$ 上次数为 $d>0$，故非空。因此，
$1:\mathcal{O}_T\to\mathcal{A}$ 可局部用作第一个基元素（例如可用
《代数》引理 \ref{algebra-lemma-cokernel-flat} 的 (1)、(2) 看出这一点）。
所以在 $T$ 上局部化后，可以假设存在同构
$\varphi : \mathcal{A} \to \mathcal{O}_T^{\oplus d}$
，使 $1\in\Gamma(\mathcal{A})$ 对应于第一个基元素。在这种情形下，
乘法映射
$\mathcal{A} \otimes_{\mathcal{O}_T} \mathcal{A} \to \mathcal{A}$
经由 $\varphi$ 转化为 $\Gamma(T,\mathcal{O}_T)$ 上的自由 $d$ 维代数
$m$。证明完成。
\end{proof}




\section{空间的有限 Hilbert 叠}
\label{criteria-section-spaces-hilbert}

\noindent
代数空间的有限 Hilbert 叠是代数叠。

\begin{lemma}
\label{criteria-lemma-hilbert-stack-of-space}
设 $S$ 是概形，$X$ 是 $S$ 上的代数空间。则 $\mathcal{H}_d(X)$ 是代数叠。
\end{lemma}

\begin{proof}
$1$-态射
$$
\mathcal{H}_d(X) \longrightarrow \mathcal{H}_d
$$
由引理 \ref{criteria-lemma-representable-on-top} 可由代数空间表示。
由命题 \ref{criteria-proposition-finite-hilbert-point}，叠 $\mathcal{H}_d$ 是代数叠。
因此，由《代数叠》引理
\ref{algebraic-lemma-representable-morphism-to-algebraic}，
$\mathcal{H}_d(X)$ 是代数叠。
\end{proof}

\noindent
此引理使我们能够进行自举。

\begin{lemma}
\label{criteria-lemma-hilbert-stack-relative-space}
设 $S$ 是概形，$F:\mathcal{X}\to\mathcal{Y}$ 是
$(\Sch/S)_{fppf}$ 上群胚叠的 $1$-态射，并满足
\begin{enumerate}
\item $\mathcal{X}$ 可由代数空间表示；
\item $F$ 可由代数空间表示、为满、平坦且局部有限表示。
\end{enumerate}
则 $\mathcal{H}_d(\mathcal{X}/\mathcal{Y})$ 是代数叠。
\end{lemma}

\begin{proof}
选取 $S$ 上的可表示群胚叠 $\mathcal{U}$ 和 $1$-态射
$f:\mathcal{U}\to\mathcal{H}_d(\mathcal{X})$，它可由代数空间表示、
光滑且为满。由于引理 \ref{criteria-lemma-hilbert-stack-of-space} 表明
$\mathcal{H}_d(\mathcal{X})$ 是代数叠，这种选择是可能的。考虑 $2$-纤维积
$$
\mathcal{W} =
\mathcal{H}_d(\mathcal{X}/\mathcal{Y})
\times_{\mathcal{H}_d(\mathcal{X}), f}
\mathcal{U}.
$$
由于 $\mathcal{U}$ 可表示（特别地，它是集合胚叠），由《叠的例子》引理
\ref{examples-stacks-lemma-faithful-hilbert} 和《叠》引理
\ref{stacks-lemma-2-fibre-product-gives-stack-in-setoids}，
$\mathcal{W}$ 是集合胚叠。$1$-态射
$\mathcal{W}\to\mathcal{H}_d(\mathcal{X}/\mathcal{Y})$ 是态射 $f$ 的基变换，
故可由代数空间表示、光滑且为满（参见《代数叠》引理
\ref{algebraic-lemma-base-change-representable-by-spaces} 和
\ref{algebraic-lemma-base-change-representable-transformations-property}）。
因此，如果能证明 $\mathcal{W}$ 可由代数空间表示，则由《代数叠》引理
\ref{algebraic-lemma-smooth-surjective-morphism-implies-algebraic}，
引理即得证。

\medskip\noindent
由引理 \ref{criteria-lemma-flat-finite-presentation-surjective-diagonal}，
$\mathcal{Y}$ 的对角态射可由代数空间表示。应用引理
\ref{criteria-lemma-relative-hilbert} 可知 $1$-态射
$$
\mathcal{H}_d(\mathcal{X}/\mathcal{Y})
\longrightarrow
\mathcal{H}_d(\mathcal{X}) \times \mathcal{Y}
$$
可由代数空间表示。考虑 $2$-纤维积
$$
\mathcal{V} =
\mathcal{H}_d(\mathcal{X}/\mathcal{Y})
\times_{(\mathcal{H}_d(\mathcal{X}) \times \mathcal{Y}), f \times F}
(\mathcal{U} \times \mathcal{X}).
$$
投影态射 $\mathcal{V}\to\mathcal{U}\times\mathcal{X}$ 是上一个所示态射的
基变换，故可由代数空间表示。因此 $\mathcal{V}$ 是代数空间（参见《自举》引理
\ref{bootstrap-lemma-representable-by-spaces-over-space}，或《代数叠》引理
\ref{algebraic-lemma-base-change-by-space-representable-by-space}）。
% P11-E0192: frozen source says V -> U here, but the displayed vertical morphism and subsequent argument are V -> W; preserved canonically in the formula.
$1$-态射 $\mathcal{V}\to\mathcal{U}$ 嵌入下列 $2$-笛卡尔图
$$
\xymatrix{
\mathcal{V} \ar[d] \ar[r] & \mathcal{X} \ar[d]^F \\
\mathcal{W} \ar[r] & \mathcal{Y}
}
$$
，这是因为
$$
\mathcal{H}_d(\mathcal{X}/\mathcal{Y})
\times_{(\mathcal{H}_d(\mathcal{X}) \times \mathcal{Y}), f \times F}
(\mathcal{U} \times \mathcal{X})
=
(\mathcal{H}_d(\mathcal{X}/\mathcal{Y})
\times_{\mathcal{H}_d(\mathcal{X}), f}
\mathcal{U}) \times_{\mathcal{Y}, F} \mathcal{X}.
$$
因此，$\mathcal{V}\to\mathcal{W}$ 作为 $F$ 的基变换，可由代数空间表示、
为满、平坦且局部有限表示。于是，与 $\mathcal{V}$ 和 $\mathcal{W}$
相伴的相应集合层也具有相同性质；参见《代数叠》引理
\ref{algebraic-lemma-map-fibred-setoids-property}。故由《自举》定理
\ref{bootstrap-theorem-final-bootstrap}，与 $\mathcal{W}$ 相伴的层
是代数空间。
\end{proof}




\section{Hilbert 叠中的 LCI 轨迹}
\label{criteria-section-lci}

\noindent
记号参见《叠的例子》第 \ref{examples-stacks-section-hilbert-d-stack} 节。
固定 $(\Sch/S)_{fppf}$ 上群胚叠的 $1$-态射
$F:\mathcal{X}\longrightarrow\mathcal{Y}$。假设 $F$ 可由代数空间表示。
固定 $d\geq1$。考虑 $\mathcal{H}_d(\mathcal{X}/\mathcal{Y})$ 的对象
$(U,Z,y,x,\alpha)$。存在诱导的 $1$-态射
$$
(\Sch/Z)_{fppf}
\longrightarrow
(\Sch/U)_{fppf} \times_{y, \mathcal{Y}, F} \mathcal{X}
$$
（由 $2$-纤维积的泛性质），它可由 $U$ 上代数空间的态射表示。
确切地，由于 $F$ 可由代数空间表示，可以选取 $U$ 上的代数空间 $X_y$，
它表示 $2$-纤维积
$(\Sch/U)_{fppf} \times_{y, \mathcal{Y}, F} \mathcal{X}$.
由于 $\alpha:y|_Z\to F(x)$ 是同构，可知
$\xi=(Z,Z\to U,x,\alpha)$ 是 $2$-纤维积
$(\Sch/U)_{fppf}\times_{y,\mathcal{Y},F}\mathcal{X}$ 在 $Z$ 上的对象。
因此，$\xi$ 给出 $U$ 上代数空间的态射 $x_\alpha:Z\to X_y$，因为
% P11-E0193: frozen source says "isomorphisms classes"; translated as the intended "isomorphism classes".
$X_y$ 是
$(\Sch/U)_{fppf} \times_{y, \mathcal{Y}, F} \mathcal{X}$
的对象同构类函子；参见《代数叠》引理
\ref{algebraic-lemma-characterize-representable-by-space}。图示为
\begin{equation}
\label{criteria-equation-relative-map}
\vcenter{
\xymatrix{
Z \ar[r]_{x_\alpha} \ar[rd] & X_y \ar[d] \\
& U
}
}
\quad\quad
\vcenter{
\xymatrix{
(\Sch/Z)_{fppf} \ar[rd] \ar[r]_-{x, \alpha} &
(\Sch/U)_{fppf} \times_{y, \mathcal{Y}, F} \mathcal{X} \ar[r] \ar[d] &
\mathcal{X} \ar[d]^F \\
& (\Sch/U)_{fppf} \ar[r]^y & \mathcal{Y}
}
}
\end{equation}
注意，如果
$(f, g, b, a) : (U, Z, y, x, \alpha) \to (U', Z', y', x', \alpha')$
是 $\mathcal{H}_d$ 的对象之间的态射，则态射
$x'_{\alpha'}:Z'\to X'_{y'}$ 是态射 $x_\alpha$ 沿 $g:U'\to U$
的基变换（略去细节）。

\medskip\noindent
现在进一步假设 $F$ 平坦且局部有限表示。在这种情形下，定义全子范畴
$$
\mathcal{H}_{d, lci}(\mathcal{X}/\mathcal{Y}) \subset
\mathcal{H}_d(\mathcal{X}/\mathcal{Y})
$$
，它由 $\mathcal{H}_d(\mathcal{X}/\mathcal{Y})$ 中满足下列条件的对象
$(U,Z,y,x,\alpha)$ 组成：相应态射 $x_\alpha:Z\to X_y$ 非分歧，
且为局部完全交态射（定义参见《空间的态射》定义
\ref{spaces-morphisms-definition-unramified} 和《空间态射进阶》定义
\ref{spaces-more-morphisms-definition-lci}）。

\begin{lemma}
\label{criteria-lemma-lci-locus-stack-in-groupoids}
设 $S$ 是概形。固定 $(\Sch/S)_{fppf}$ 上群胚叠的 $1$-态射
$F:\mathcal{X}\longrightarrow\mathcal{Y}$。假设 $F$ 可由代数空间表示、
平坦且局部有限表示。则 $\mathcal{H}_{d,lci}(\mathcal{X}/\mathcal{Y})$
是群胚叠，且嵌入函子
$$
\mathcal{H}_{d, lci}(\mathcal{X}/\mathcal{Y})
\longrightarrow
\mathcal{H}_d(\mathcal{X}/\mathcal{Y})
$$
可表示且为开浸入。
\end{lemma}

\begin{proof}
设 $\Xi=(U,Z,y,x,\alpha)$ 是 $\mathcal{H}_d$ 的对象。由
(\ref{criteria-equation-relative-map}) 后的评注，$\Xi$ 沿 $U'\to U$ 的拉回属于
$\mathcal{H}_{d,lci}(\mathcal{X}/\mathcal{Y})$，当且仅当 $x_\alpha$
的基变换非分歧且为局部完全交态射。注意，$Z\to U$ 有限局部自由
（因而平坦、局部有限表示且万有闭），而由关于 $F$ 的假设，$X_y\to U$
平坦且局部有限表示。于是《空间态射进阶》引理
\ref{spaces-more-morphisms-lemma-where-unramified} 和
\ref{spaces-more-morphisms-lemma-where-lci}
% P11-E0194: frozen source says "imply exists"; "there" is missing. Translated naturally.
表明存在开子概形 $W\subset U$，使态射 $U'\to U$ 通过 $W$ 分解，
当且仅当 $x_\alpha$ 沿 $U'\to U$ 的基变换非分歧且为局部完全交态射。
这意味着
$$
(\Sch/U)_{fppf}
\times_{\Xi, \mathcal{H}_d(\mathcal{X}/\mathcal{Y})}
\mathcal{H}_{d, lci}(\mathcal{X}/\mathcal{Y})
$$
可由 $W$ 表示。因此，引理的最后一个断言成立。第一个断言（即
$\mathcal{H}_{d,lci}(\mathcal{X}/\mathcal{Y})$ 是群胚叠）由此和
《代数叠》引理 \ref{algebraic-lemma-open-fibred-category-is-algebraic}
推出。
\end{proof}

\noindent
局部完全交态射“局部无阻碍”。这一点在远比本章所需特殊情形更一般的
范围内成立。

\begin{lemma}
\label{criteria-lemma-lci-unobstructed}
设 $U\subset U'$ 是仿射概形的一阶加厚，$X'$ 是 $U'$ 上平坦的代数空间。
令 $X=U\times_{U'}X'$。设 $Z\to U$ 是次数为 $d$ 的有限局部自由态射。
最后，设 $f:Z\to X$ 非分歧且为局部完全交态射。则存在交换图
$$
\xymatrix{
(Z \subset Z') \ar[rd] \ar[rr]_{(f, f')} & & (X \subset X') \ar[ld] \\
& (U \subset U')
}
$$
，其中各项为 $U'$ 上的代数空间，且 $Z'\to U'$ 是次数为 $d$ 的有限局部自由
态射，并满足 $Z=U\times_{U'}Z'$。
\end{lemma}

\begin{proof}
由《空间态射进阶》引理 \ref{spaces-more-morphisms-lemma-unramified-lci}，
非分歧态射 $Z\to X$ 的余法层 $\mathcal{C}_{Z/X}$ 是有限局部自由
$\mathcal{O}_Z$-模；又由《空间态射进阶》引理
\ref{spaces-more-morphisms-lemma-transitivity-conormal-lci}，有余法层的正合列
$$
0 \to i^*\mathcal{C}_{X/X'} \to
\mathcal{C}_{Z/X'} \to
\mathcal{C}_{Z/X} \to 0
$$
。由于 $Z$ 仿射，此列分裂。选取分裂
$$
\mathcal{C}_{Z/X'} = i^*\mathcal{C}_{X/X'} \oplus \mathcal{C}_{Z/X}
$$
设 $Z\subset Z''$ 是 $Z$ 在 $X'$ 上的万有一阶加厚（参见《空间态射进阶》第
\ref{spaces-more-morphisms-section-universal-thickening} 节）。
% P11-E0195: frozen source's naming construction "Denote I ... the" omits "by"; translated naturally.
记 $\mathcal{I}\subset\mathcal{O}_{Z''}$ 为与 $Z\subset Z''$ 相应的拟凝聚理想层。
% P11-E0196: frozen source says "we have C is I"; the intended identification is translated naturally.
按定义，把 $\mathcal{I}$ 视为 $Z$ 上的层时，它就是 $\mathcal{C}_{Z/X'}$。
因此，上述分裂确定分裂
$$
\mathcal{I} = i^*\mathcal{C}_{X/X'} \oplus \mathcal{C}_{Z/X}
$$
把 $\mathcal{C}_{Z/X}\subset\mathcal{I}$ 视为 $Z''$ 上的拟凝聚理想层，
令 $Z'\subset Z''$ 为它截出的闭子概形。显然，$Z'$ 是 $Z$ 的一阶加厚，
并得到引理陈述中的一阶加厚交换图。

\medskip\noindent
由于 $X'\to U'$ 平坦且 $X=U\times_{U'}X'$，可知
$\mathcal{C}_{X/X'}$ 是 $\mathcal{C}_{U/U'}$ 到 $X$ 的拉回；参见
《空间态射进阶》引理 \ref{spaces-more-morphisms-lemma-deform}。
注意，按构造，$\mathcal{C}_{Z/Z'}=i^*\mathcal{C}_{X/X'}$，故
$\mathcal{C}_{Z/Z'}$ 同构于 $\mathcal{C}_{U/U'}$ 到 $Z$ 的拉回。
再次应用《空间态射进阶》引理 \ref{spaces-more-morphisms-lemma-deform}
（或其概形版本，参见《态射进阶》引理 \ref{more-morphisms-lemma-deform}），
可知 $Z'\to U'$ 平坦且 $Z=U\times_{U'}Z'$。最后，《态射进阶》引理
\ref{more-morphisms-lemma-deform-property} 表明 $Z'\to U'$ 是次数为 $d$
的有限局部自由态射。
\end{proof}

\begin{lemma}
\label{criteria-lemma-lci-formally-smooth}
设 $F:\mathcal{X}\to\mathcal{Y}$ 是 $(\Sch/S)_{fppf}$ 上群胚叠的
$1$-态射。假设 $F$ 可由代数空间表示、平坦且局部有限表示。则
$$
p : \mathcal{H}_{d, lci}(\mathcal{X}/\mathcal{Y}) \to \mathcal{Y}
$$
在对象上形式光滑。
\end{lemma}

\begin{proof}
须证明如下命题：给定
\begin{enumerate}
\item $\mathcal{H}_{d,lci}(\mathcal{X}/\mathcal{Y})$ 在仿射概形 $U$
上的对象 $(U,Z,y,x,\alpha)$；
\item 一阶加厚 $U\subset U'$；以及
\item $\mathcal{Y}$ 在 $U'$ 上的对象 $y'$，满足 $y'|_U=y$，
\end{enumerate}
则存在 $\mathcal{H}_{d,lci}(\mathcal{X}/\mathcal{Y})$ 在 $U'$ 上的对象
$(U',Z',y',x',\alpha')$，满足 $Z=U\times_{U'}Z'$、$x=x'|_Z$ 以及
$\alpha=\alpha'|_U$。最后两个等式恰好保证
(\ref{criteria-equation-formally-smooth}) 交换。

\medskip\noindent
考虑公式 (\ref{criteria-equation-relative-map}) 中构造的态射
$x_\alpha:Z\to X_y$。类似地，
% P11-E0197: frozen source's naming construction "Denote similarly X' the ..." omits "by"; translated naturally.
记 $X'_{y'}$ 为 $U'$ 上表示 $2$-纤维积
$(\Sch/U')_{fppf} \times_{y', \mathcal{Y}, F} \mathcal{X}$.
按假设，态射 $X'_{y'}\to U'$ 平坦（且局部有限表示）。由于
$y'|_U=y$，可知 $X_y=U\times_{U'}X'_{y'}$。因此，可以应用引理
\ref{criteria-lemma-lci-unobstructed}，找到次数为 $d$ 的有限局部自由态射
$Z'\to U'$，满足 $Z=U\times_{U'}Z'$，并使 $Z'\to X'_{y'}$ 延拓
$x_\alpha$。按构造，态射 $Z'\to X'_{y'}$ 对应于二元组 $(x',\alpha')$。
显然，$(U',Z',y',x',\alpha')$ 是
$\mathcal{H}_d(\mathcal{X}/\mathcal{Y})$ 在 $U'$ 上的对象，并满足
$Z=U\times_{U'}Z'$、$x=x'|_Z$ 和 $\alpha=\alpha'|_U$。
由引理 \ref{criteria-lemma-lci-locus-stack-in-groupoids} 已知
$\mathcal{H}_{d,lci}(\mathcal{X}/\mathcal{Y})\subset
\mathcal{H}_d(\mathcal{X}/\mathcal{Y})$ 是“开子叠”，故
$(U',Z',y',x',\alpha')$ 如愿是
$\mathcal{H}_{d,lci}(\mathcal{X}/\mathcal{Y})$ 的对象。
\end{proof}

\begin{lemma}
\label{criteria-lemma-lci-surjective}
设 $F:\mathcal{X}\to\mathcal{Y}$ 是 $(\Sch/S)_{fppf}$ 上群胚叠的
$1$-态射。假设 $F$ 可由代数空间表示、平坦、为满且局部有限表示。则
$$
\coprod\nolimits_{d \geq 1} \mathcal{H}_{d, lci}(\mathcal{X}/\mathcal{Y})
\longrightarrow
\mathcal{Y}
$$
在对象上为满。
\end{lemma}

\begin{proof}
只需证明如下命题：对任意域 $k$ 和 $\mathcal{Y}$ 在 $\Spec(k)$ 上的对象
$y$，存在整数 $d\geq1$ 和
$\mathcal{H}_{d,lci}(\mathcal{X}/\mathcal{Y})$ 的对象
$(U,Z,y,x,\alpha)$，其中 $U=\Spec(k)$。此时，$p$ 以无需域扩张的
强意义在对象上为满。

\medskip\noindent
% P11-E0198: frozen source's naming construction "Denote X_y the ..." omits "by"; translated naturally.
记 $X_y$ 为 $U=\Spec(k)$ 上表示 $2$-纤维积
% P11-E0199: frozen source uses U' and y' in the following fibre product although this proof has introduced U and y only; frozen math is preserved.
$(\Sch/U')_{fppf} \times_{y', \mathcal{Y}, F} \mathcal{X}$.
的代数空间。按假设，态射 $X_y\to\Spec(k)$ 为满且局部有限表示
（并且平坦）。特别地，$X_y$ 非空。选取非空仿射概形 $V$ 和平展态射
$V\to X_y$。注意，$V\to\Spec(k)$（平坦）、为满且局部有限表示
（由《空间的态射》定义
\ref{spaces-morphisms-definition-locally-finite-presentation}）。
选取闭点 $v\in V$，使 $V\to\Spec(k)$ 在该点为 Cohen--Macaulay
（即 $V$ 在 $v$ 处为 Cohen--Macaulay）；参见《态射进阶》引理
\ref{more-morphisms-lemma-flat-finite-presentation-CM-open}。应用
《态射进阶》引理 \ref{more-morphisms-lemma-slice-CM}，得到正则浸入
$Z\to V$，其中 $Z=\{v\}$。这蕴含 $Z\to V$ 是闭浸入。此外，
$Z\to\Spec(k)$ 有限（例如由《代数》引理
\ref{algebra-lemma-isolated-point}）。故 $Z\to\Spec(k)$ 是某个次数 $d$
的有限局部自由态射。现在，$Z\to X_y$ 是闭浸入后接平展态射的复合，
因而非分歧（参见《空间的态射》引理
\ref{spaces-morphisms-lemma-composition-unramified},
\ref{spaces-morphisms-lemma-etale-unramified} 和
\ref{spaces-morphisms-lemma-immersion-unramified}）。
最后，$Z \to X_y$ 是局部完全交态射，因为它是概形的正则浸入与
代数空间的平展态射的复合（参见《态射进阶》引理
\ref{more-morphisms-lemma-regular-immersion-lci}，
《空间的态射》引理 \ref{spaces-morphisms-lemma-etale-smooth} 和
\ref{spaces-morphisms-lemma-smooth-syntomic}，以及
《空间的态射进阶》引理
\ref{spaces-more-morphisms-lemma-flat-lci} 和
\ref{spaces-more-morphisms-lemma-composition-lci}）。
态射 $Z \to X_y$ 对应于 $\mathcal{X}$ 在 $Z$ 上的一个对象 $x$，
以及同构 $\alpha : y|_Z \to F(x)$。我们由此得到
$\mathcal{H}_d(\mathcal{X}/\mathcal{Y})$ 的一个对象
$(U, Z, y, x, \alpha)$。由上面对态射 $Z \to X_y$ 的论述可知，
它实际上是子范畴
$\mathcal{H}_{d, lci}(\mathcal{X}/\mathcal{Y})$ 的对象，结论得证。
\end{proof}




















\section{代数叠的自举}
\label{criteria-section-bootstrap}

\noindent
下述定理是本章的主要结果之一。

\begin{theorem}
\label{criteria-theorem-bootstrap}
\begin{slogan}
Artin 关于平坦群胚可表示性的定理
\end{slogan}
设 $S$ 是概形，$F : \mathcal{X} \to \mathcal{Y}$ 是
$(\Sch/S)_{fppf}$ 上群胚叠的一个 $1$-态射。若
\begin{enumerate}
\item $\mathcal{X}$ 可由代数空间表示，并且
\item $F$ 可由代数空间表示、为满、平坦且局部有限表示，
\end{enumerate}
则 $\mathcal{Y}$ 是代数叠。
\end{theorem}

\begin{proof}
由引理 \ref{criteria-lemma-flat-finite-presentation-surjective-diagonal}，
$\mathcal{Y}$ 的对角可由代数空间表示。因此，我们只需验证存在
$(\Sch/S)_{fppf}$ 上群胚叠的一个 $1$-态射
$f : \mathcal{V} \to \mathcal{Y}$，使得 $\mathcal{V}$ 可表示，
而 $f$ 为满且光滑。由引理 \ref{criteria-lemma-hilbert-stack-relative-space}，
我们知道
$$
\coprod\nolimits_{d \geq 1} \mathcal{H}_d(\mathcal{X}/\mathcal{Y})
$$
是代数叠。由引理 \ref{criteria-lemma-lci-locus-stack-in-groupoids} 和
《代数叠》引理 \ref{algebraic-lemma-open-fibred-category-is-algebraic}，
可知
$$
\coprod\nolimits_{d \geq 1} \mathcal{H}_{d, lci}(\mathcal{X}/\mathcal{Y})
$$
也是代数叠。选取 $(\Sch/S)_{fppf}$ 上一个可表示的群胚叠
$\mathcal{V}$，以及一个满且光滑的 $1$-态射
$$
\mathcal{V}
\longrightarrow
\coprod\nolimits_{d \geq 1} \mathcal{H}_{d, lci}(\mathcal{X}/\mathcal{Y}).
$$
我们断言复合态射
$$
\mathcal{V}
\longrightarrow
\coprod\nolimits_{d \geq 1} \mathcal{H}_{d, lci}(\mathcal{X}/\mathcal{Y})
\longrightarrow
\mathcal{Y}
$$
光滑且为满，这就完成定理的证明。事实上，光滑性将由引理
\ref{criteria-lemma-limit-preserving} 和 \ref{criteria-lemma-lci-formally-smooth}
推出，而满性将由引理 \ref{criteria-lemma-lci-surjective} 推出。
我们在下一段详述其细节。

\medskip\noindent
由构造，$\mathcal{V} \to
\coprod\nolimits_{d \geq 1} \mathcal{H}_{d, lci}(\mathcal{X}/\mathcal{Y})$
可由代数空间表示、为满且光滑（因而由一般原理
《代数叠》引理
\ref{algebraic-lemma-representable-transformations-property-implication}
和《空间的态射进阶》引理
\ref{spaces-more-morphisms-lemma-smooth-formally-smooth}，
它也局部有限表示且形式光滑）。应用引理
\ref{criteria-lemma-representable-by-spaces-limit-preserving}、
\ref{criteria-lemma-representable-by-spaces-formally-smooth} 和
\ref{criteria-lemma-representable-by-spaces-surjective}
可知 $\mathcal{V} \to
\coprod\nolimits_{d \geq 1} \mathcal{H}_{d, lci}(\mathcal{X}/\mathcal{Y})$
在对象上保极限、在对象上形式光滑且在对象上为满。$1$-态射
$\coprod\nolimits_{d \geq 1} \mathcal{H}_{d, lci}(\mathcal{X}/\mathcal{Y})
\to \mathcal{Y}$ 满足：
\begin{enumerate}
\item 在对象上保极限：对
$\mathcal{H}_d(\mathcal{X}/\mathcal{Y}) \to \mathcal{Y}$ 而言，
这是引理 \ref{criteria-lemma-limit-preserving}；将其与引理
\ref{criteria-lemma-lci-locus-stack-in-groupoids}、
\ref{criteria-lemma-open-immersion-limit-preserving} 和
\ref{criteria-lemma-composition-limit-preserving}
结合，便得到
$\mathcal{H}_{d, lci}(\mathcal{X}/\mathcal{Y}) \to \mathcal{Y}$
的相应结论；
\item 由引理 \ref{criteria-lemma-lci-formally-smooth}，在对象上形式光滑；并且
\item 由引理 \ref{criteria-lemma-lci-surjective}，在对象上为满。
\end{enumerate}
应用引理 \ref{criteria-lemma-composition-limit-preserving}、
\ref{criteria-lemma-composition-formally-smooth} 和
\ref{criteria-lemma-composition-surjective}
可知复合态射 $\mathcal{V} \to \mathcal{Y}$ 在对象上保极限、
在对象上形式光滑且在对象上为满。再应用引理
\ref{criteria-lemma-representable-by-spaces-limit-preserving}、
\ref{criteria-lemma-representable-by-spaces-formally-smooth} 和
\ref{criteria-lemma-representable-by-spaces-surjective}
可知 $\mathcal{V} \to \mathcal{Y}$ 局部有限表示、形式光滑且为满。
最后，通过一般原理《代数叠》引理
\ref{algebraic-lemma-representable-transformations-property-implication}
应用无穷小提升判据（《空间的态射进阶》引理
\ref{spaces-more-morphisms-lemma-smooth-formally-smooth}），可知
$\mathcal{V} \to \mathcal{Y}$ 光滑，结论得证。
\end{proof}








\section{应用}
\label{criteria-section-applications}

\noindent
我们的首要任务是证明，与“平坦且局部有限表示的群胚”相伴的商叠
$[U/R]$ 是代数叠。商叠的定义见《空间中的群胚》定义
\ref{spaces-groupoids-definition-quotient-stack}。下述预备引理是
《代数叠》引理
\ref{algebraic-lemma-smooth-quotient-smooth-presentation} 的对应版本。

\begin{lemma}
\label{criteria-lemma-flat-quotient-flat-presentation}
设 $S$ 是 $\Sch_{fppf}$ 中的概形，
$(U, R, s, t, c)$ 是 $S$ 上代数空间中的群胚。
假设 $s, t$ 平坦且局部有限表示。则态射
$\mathcal{S}_U \to [U/R]$ 平坦、局部有限表示且为满。
\end{lemma}

\begin{proof}
设 $T$ 是概形，$x : (\Sch/T)_{fppf} \to [U/R]$ 是一个
$1$-态射。我们必须证明投影
$$
\mathcal{S}_U \times_{[U/R]} (\Sch/T)_{fppf}
\longrightarrow
(\Sch/T)_{fppf}
$$
为满、平坦且局部有限表示。我们已经知道左边可由一个代数空间 $F$
表示，参见《代数叠》引理
\ref{algebraic-lemma-diagonal-quotient-stack} 和
\ref{algebraic-lemma-representable-diagonal}。因此，我们必须证明相应的
代数空间态射 $F \to T$ 为满、局部有限表示且平坦。由于此处所考察的
代数空间态射性质关于 fppf 拓扑在靶上是局部的，我们可以在 $T$ 上
fppf 局部地检验这些性质。由构造，存在 $T$ 的一个 fppf 覆盖
$\{T_i \to T\}$，使得 $x|_{(\Sch/T_i)_{fppf}}$ 来自某个态射
$x_i : T_i \to U$。（注意，$F \times_T T_i$ 表示 $2$-纤维积
$\mathcal{S}_U \times_{[U/R]} (\Sch/T_i)_{fppf}$，故一切都与经由
$T_i \to T$ 的基变换相容。）因此，可假设 $x$ 来自
$x : T \to U$。在这种情形下，我们看到
$$
\mathcal{S}_U \times_{[U/R]} (\Sch/T)_{fppf}
=
(\mathcal{S}_U \times_{[U/R]} \mathcal{S}_U)
\times_{\mathcal{S}_U} (\Sch/T)_{fppf}
=
\mathcal{S}_R \times_{\mathcal{S}_U} (\Sch/T)_{fppf}
$$
其中第一个等号由《范畴》引理
\ref{categories-lemma-2-fibre-product-erase-factor} 给出，第二个等号由
《空间中的群胚》引理
\ref{spaces-groupoids-lemma-quotient-stack-2-cartesian} 给出。显然，
最后一个 $2$-纤维积由代数空间 $F = R \times_{s, U, x} T$ 表示；
投影 $R \times_{s, U, x} T \to T$ 是平坦且局部有限表示的代数空间态射
$s : R \to U$ 的基变换，因而平坦且局部有限表示。它也为满，因为
$s$ 有截面（即群胚的恒等态射 $e : U \to R$）。引理得证。
\end{proof}

\noindent
下面是本节的第一个主要结果。

\begin{theorem}
\label{criteria-theorem-flat-groupoid-gives-algebraic-stack}
设 $S$ 是 $\Sch_{fppf}$ 中的概形，
$(U, R, s, t, c)$ 是 $S$ 上代数空间中的群胚。
假设 $s, t$ 平坦且局部有限表示。则商叠 $[U/R]$ 是 $S$ 上的代数叠。
\end{theorem}

\begin{proof}
对态射
$$
(\Sch/U)_{fppf} \longrightarrow [U/R].
$$
我们检验定理 \ref{criteria-theorem-bootstrap} 的两个条件。第一个条件是显然的
（因为 $U$ 是代数空间）。第二个条件即引理
\ref{criteria-lemma-flat-quotient-flat-presentation}。
\end{proof}









\section{商叠何时是代数的？}
\label{criteria-section-quotient-algebraic}

\noindent
在《空间中的群胚》第 \ref{spaces-groupoids-section-stacks} 节中，
我们定义了与代数空间中的群胚 $(U, R, s, t, c)$ 相伴的商叠 $[U/R]$。
注意，$[U/R]$ 是一个群胚叠，其对角可由代数空间表示（参见
《自举》引理 \ref{bootstrap-lemma-quotient-stack-isom} 和
《代数叠》引理 \ref{algebraic-lemma-representable-diagonal}），并且存在
代数空间 $U$ 和一个 $1$-态射 $(\Sch/U)_{fppf} \to [U/R]$，它是如下
意义下的“fppf 满射”：它在对象同构类的预层上诱导的映射经层化后为满。
然而，一般而言 $[U/R]$ 并不一定是代数叠。这与定理
\ref{criteria-theorem-bootstrap} 并不矛盾，因为 $1$-态射
$(\Sch/U)_{fppf} \to [U/R]$ 未必平坦且局部有限表示。

\medskip\noindent
构造非代数商叠例子的最简方法，是考察形如 $[S/G]$ 的商，其中 $S$
是概形，而 $G$ 是 $S$ 上平凡作用于 $S$ 的群概形。确切地说，下面将看到
（引理 \ref{criteria-lemma-BG-algebraic}），若 $[S/G]$ 是代数的，则
$G \to S$ 必须平坦且局部有限表示。一个显式例子见《例子》第
\ref{examples-section-not-algebraic-stack} 节。

\begin{lemma}
\label{criteria-lemma-quotient-algebraic}
设 $S$ 是概形，$B$ 是 $S$ 上的代数空间，
$(U, R, s, t, c)$ 是 $B$ 上代数空间中的群胚。商叠 $[U/R]$
是代数叠，当且仅当存在代数空间态射 $g : U' \to U$，使得
\begin{enumerate}
\item 复合态射
$U' \times_{g, U, t} R \to R \xrightarrow{s} U$ 是层的满射；并且
\item 态射 $s', t' : R' \to U'$ 平坦且局部有限表示，其中
$(U', R', s', t', c')$ 是 $(U, R, s, t, c)$ 沿 $g$ 的限制。
\end{enumerate}
\end{lemma}

\begin{proof}
首先，假设 $g : U' \to U$ 满足 (1) 和 (2)。性质 (1) 蕴含
$[U'/R'] \to [U/R]$ 是等价，参见《空间中的群胚》引理
\ref{spaces-groupoids-lemma-quotient-stack-restrict-equivalence}。
由定理 \ref{criteria-theorem-flat-groupoid-gives-algebraic-stack}，商叠
$[U'/R']$ 是代数叠。因此 $[U/R]$ 也是代数叠，参见《代数叠》引理
\ref{algebraic-lemma-equivalent}。

\medskip\noindent
反过来，假设 $[U/R]$ 是代数叠。我们可以选取概形 $W$ 和一个满光滑
$1$-态射
$$
f : (\Sch/W)_{fppf} \longrightarrow [U/R].
$$
由 $2$-Yoneda 引理（《代数叠》第 \ref{algebraic-section-2-yoneda} 节），
这对应于 $[U/R]$ 在 $W$ 上的一个对象 $\xi$。由《空间中的群胚》引理
\ref{spaces-groupoids-lemma-quotient-stack-objects} 对 $[U/R]$ 的描述，
我们可以找到概形的一个满、平坦且局部有限表示的态射
$b : U' \to W$，使得 $\xi' = b^*\xi$ 对应于某个态射
$g : U' \to U$。注意，与 $\xi'$ 对应的 $1$-态射
$$
f' : (\Sch/U')_{fppf} \longrightarrow [U/R].
$$
为满、平坦且局部有限表示，参见《代数叠》引理
\ref{algebraic-lemma-composition-representable-transformations-property}。
因此，
$(\Sch/U')_{fppf} \times_{[U/R]} (\Sch/U')_{fppf}$
由代数空间
$$
\mathit{Isom}_{[U/R]}(\text{pr}_0^*\xi', \text{pr}_1^*\xi') =
(U' \times_S U')
\times_{(g \circ \text{pr}_0, g \circ \text{pr}_1), U \times_S U} R = R'
$$
表示（第一个等号见《空间中的群胚》引理
\ref{spaces-groupoids-lemma-quotient-stack-morphisms}；第二个等号是限制的
定义），并且经由 $s'$ 和 $t'$ 在 $U'$ 上都平坦且局部有限表示
（由基变换，参见《代数叠》引理
\ref{algebraic-lemma-base-change-representable-transformations-property}）。
由这一对 $R'$ 的描述以及《代数叠》引理
\ref{algebraic-lemma-map-space-into-stack}，我们得到一个典范的全忠实
$1$-态射 $[U'/R'] \to [U/R]$。因为 $f'$ 平坦、局部有限表示且为满，
这个 $1$-态射本质满（参见《叠》引理
\ref{stacks-lemma-characterize-essentially-surjective-when-ff}）；另一种证明
方法是应用《代数叠》注记
\ref{algebraic-remark-flat-fp-presentation}。最后，应用《空间中的群胚》引理
\ref{spaces-groupoids-lemma-quotient-stack-restrict-equivalence}，可知复合态射
$U' \times_{g, U, t} R \to R \xrightarrow{s} U$ 是层的满射。
\end{proof}

\begin{lemma}
\label{criteria-lemma-group-quotient-algebraic}
设 $S$ 是概形，$B$ 是 $S$ 上的代数空间，$G$ 是 $B$ 上的群代数空间。
设 $X$ 是 $B$ 上的代数空间，而 $a : G \times_B X \to X$ 是 $G$
在 $B$ 上对 $X$ 的一个作用。商叠 $[X/G]$ 是代数叠，当且仅当存在
代数空间态射 $\varphi : X' \to X$，使得
\begin{enumerate}
\item $G \times_B X' \to X$，$(g, x') \mapsto a(g, \varphi(x'))$，
是层的满射；并且
\item 由规则
$$
T \longmapsto \{(x'_1, g, x'_2) \in (X' \times_B G \times_B X')(T)
\mid \varphi(x'_1) = a(g, \varphi(x'_2))\}
$$
给出的代数空间 $X''$ 的两个投影 $X'' \to X'$ 都平坦且局部有限表示。
\end{enumerate}
\end{lemma}

\begin{proof}
本引理是引理 \ref{criteria-lemma-quotient-algebraic} 的一个特例。确切地说，
由《空间中的群胚》定义
\ref{spaces-groupoids-definition-quotient-stack}，商叠 $[X/G]$ 等于
代数空间中的群胚 $(X, G \times_B X, s, t, c)$ 的商叠
$[X/G \times_B X]$；该群胚与《空间中的群胚》引理
\ref{spaces-groupoids-lemma-groupoid-from-action} 中的群作用相伴。
要得到条件 (1)，还需作一个小观察。态射 $s : G \times_B X \to X$
是第二投影，而态射 $t :  G \times_B X \to X$ 是作用态射 $a$。
因此，引理 \ref{criteria-lemma-quotient-algebraic} 中的态射
$h : U' \times_{g, U, t} R \to R \xrightarrow{s} U$ 对应于当前设置中的态射
$$
X' \times_{\varphi, X, a} (G \times_B X) \xrightarrow{\text{pr}_1} X
$$
不过，由 $G$ 的取逆所给出的对称性，此态射同构于引理陈述中的态射
$$
(G \times_B X) \times_{\text{pr}_1, X, \varphi} X' \xrightarrow{a} X
$$
。细节从略。
\end{proof}

\begin{lemma}
\label{criteria-lemma-BG-algebraic}
\begin{slogan}
Gerbe 是代数的，当且仅当相伴的群平坦且局部有限表示
\end{slogan}
设 $S$ 是概形，$B$ 是 $S$ 上的代数空间，$G$ 是 $B$ 上的群代数空间。
赋予 $B$ 以 $G$ 的平凡作用。则商叠 $[B/G]$ 是代数叠，当且仅当
$G$ 在 $B$ 上平坦且局部有限表示。
\end{lemma}

\begin{proof}
若 $G$ 在 $B$ 上平坦且局部有限表示，则由定理
\ref{criteria-theorem-flat-groupoid-gives-algebraic-stack}，$[B/G]$ 是代数叠。

\medskip\noindent
反过来，假设 $[B/G]$ 是代数叠。由引理
\ref{criteria-lemma-group-quotient-algebraic} 以及作用的平凡性，可知存在代数空间
$B'$ 和态射 $B' \to B$，使得 (1) $B' \to B$ 是层的满射，并且
(2) 投影
$$
B' \times_B G \times_B B' \to B'
$$
平坦且局部有限表示。注意，$B' \to B$ 的基变换
$B' \times_B G \times_B B' \to G \times_B B'$ 也是层的满射。因此，
由《空间上的下降》引理 \ref{spaces-descent-lemma-curiosity}，投影
$G \times_B B' \to B'$ 平坦且局部有限表示。由 (1)，可以找到一个
fppf 覆盖 $\{B_i \to B\}$，使得 $B_i \to B$ 经 $B' \to B$ 分解。
于是，由基变换，$G \times_B B_i \to B_i$ 平坦且局部有限表示。由
《空间上的下降》引理
\ref{spaces-descent-lemma-descending-property-flat} 和
\ref{spaces-descent-lemma-descending-property-locally-finite-presentation}，
可知 $G \to B$ 平坦且局部有限表示。
\end{proof}

\noindent
稍后将看到，光滑 $S$-空间关于群代数空间 $G$ 的商叠是光滑的，
即使 $G$ 并不光滑亦然（《叠的态射》引理
\ref{stacks-morphisms-lemma-smooth-quotient-stack}）。




\section{\'etale 拓扑中的代数叠}
\label{criteria-section-stacks-etale}

\noindent
设 $S$ 是概形。我们可以研究大 fppf 位点 $(\Sch/S)_{fppf}$ 上的群胚叠，
也可以改为研究大 \'etale 位点 $(\Sch/S)_\etale$ 上的群胚叠。
《代数叠》第
\ref{algebraic-section-representable},
\ref{algebraic-section-2-yoneda},
\ref{algebraic-section-representable-morphism},
\ref{algebraic-section-split},
\ref{algebraic-section-representable-by-algebraic-spaces},
\ref{algebraic-section-morphisms-representable-by-algebraic-spaces},
\ref{algebraic-section-representable-properties} 和
\ref{algebraic-section-stacks}
节的全部内容，对 $(\Sch/S)_\etale$ 上群胚纤维化范畴同样有意义。因此，
在 \'etale 拓扑中工作会给出代数叠的第二种概念。这个概念先验地弱于
《代数叠》定义 \ref{algebraic-definition-algebraic-stack} 引入的概念，
因为 fppf 拓扑中的叠当然也是 \'etale 拓扑中的叠。不过，下述引理表明
这两个概念等价。

\begin{lemma}
\label{criteria-lemma-stacks-etale}
用 $\Sch_\alpha$ 表示 $\Sch_{fppf}$ 与 $\Sch_\etale$ 共同的底层范畴
（参见《叠上的层》第 \ref{stacks-sheaves-section-sheaves} 节和
《拓扑》注记 \ref{topologies-remark-choice-sites}）。设 $S$ 是
$\Sch_\alpha$ 的对象，并设
$$
p : \mathcal{X} \to \Sch_\alpha/S
$$
是满足下列性质的群胚纤维化范畴：
\begin{enumerate}
\item $\mathcal{X}$ 是 $(\Sch/S)_\etale$ 上的群胚叠；
\item 对角 $\Delta : \mathcal{X} \to \mathcal{X} \times \mathcal{X}$
可由代数空间表示\footnote{这里，“代数空间”既可以指 \'etale 拓扑中
对角可表示且具有概形给出的满 \'etale 覆盖的层，也可以指《代数空间》
定义 \ref{spaces-definition-algebraic-space} 所定义的代数空间。事实上，
由《自举》引理 \ref{bootstrap-lemma-spaces-etale}，两者没有区别。}；并且
\item 存在 $U \in \Ob(\Sch_\alpha/S)$ 和一个满光滑 $1$-态射
$(\Sch/U)_\etale \to \mathcal{X}$。
\end{enumerate}
则 $\mathcal{X}$ 是《代数叠》定义
\ref{algebraic-definition-algebraic-stack} 意义下的代数叠。
\end{lemma}

\begin{proof}
注意，本引理的性质 (2)、(3) 以及《代数叠》定义
\ref{algebraic-definition-algebraic-stack} 中相应的性质 (2)、(3)
与拓扑无关。这是因为这些性质只涉及群胚纤维化范畴的 $2$-纤维积概念、
群胚纤维化范畴的 $1$-态射和 $2$-态射、可由代数空间表示的群胚纤维化
范畴 $1$-态射这一概念，以及这种 $1$-态射为满且光滑的含义。因此，
我们只需证明，满足性质 (2)、(3) 的 \'etale 群胚叠 $\mathcal{X}$
也是 fppf 群胚叠。

\medskip\noindent
应用 (2)，令 $R$ 为表示下式的代数空间：
$$
(\Sch_\alpha/U) \times_\mathcal{X} (\Sch_\alpha/U)
$$
由 (3)，投影 $s, t : R \to U$ 光滑。完全仿照《代数叠》引理
\ref{algebraic-lemma-map-space-into-stack} 的证明，存在空间中的群胚
$(U, R, s, t, c)$ 和典范全忠实 $1$-态射
$[U/R]_\etale \to \mathcal{X}$，其中 $[U/R]_\etale$ 是群胚预层
$$
T \longmapsto (U(T), R(T), s(T), t(T), c(T))
$$
的 \'etale 叠化。断言：若 $V \to T$ 是从代数空间 $V$ 到概形 $T$
的满光滑态射，则存在一个加细覆盖 $\{V \to T\}$ 的 \'etale 覆盖
$\{T_i \to T\}$。这由《态射进阶》引理
\ref{more-morphisms-lemma-etale-dominates-smooth} 或更一般的《叠上的层》
引理 \ref{stacks-sheaves-lemma-surjective-flat-locally-finite-presentation}
推出。应用此断言，并完全仿照《代数叠》引理
\ref{algebraic-lemma-stack-presentation} 的论证，可知
$[U/R]_\etale \to \mathcal{X}$ 是等价。

\medskip\noindent
接着，用 $[U/R]$ 表示 fppf 拓扑中的商叠；由《代数叠》定理
\ref{algebraic-theorem-smooth-groupoid-gives-algebraic-stack}，它是代数叠。
于是有 $1$-态射
$$
U \to [U/R]_\etale \to [U/R].
$$
$U \to [U/R]_\etale \cong \mathcal{X}$ 与 $U \to [U/R]$ 都满且光滑
（前者由假设，后者由该定理），而且在这两种情形下，纤维积
$U \times_\mathcal{X} U$ 和 $U \times_{[U/R]} U$ 都由 $R$ 表示。因此，
$1$-态射 $[U/R]_\etale \to [U/R]$ 全忠实（因为商叠中的态射由到
$R$ 的态射给出，参见《空间中的群胚》第
\ref{spaces-groupoids-section-explicit-quotient-stacks} 节）。

\medskip\noindent
最后，对任意概形 $T$ 和态射 $t : T \to [U/R]$，纤维积
$V = T \times_{U/R} U$ 是在 $T$ 上满且光滑的代数空间。由上述断言，
存在一个 \'etale 覆盖 $\{T_i \to T\}_{i \in I}$ 以及 $T$ 上的态射
$T_i \to V$。这证明 $[U/R]$ 在 $T$ 上的对象 $t$ 在 \'etale 局部来自
$U$。由《叠》引理
\ref{stacks-lemma-characterize-essentially-surjective-when-ff}，可知
$[U/R]_\etale \to [U/R]$ 是 $(\Sch/S)_\etale$ 上群胚叠的等价。
证明完毕。
\end{proof}
